% NAL 1644, batch 3 — Gallica views 41–60.
% Pass: Opus 5.5 (claude-opus-5-5), 2026-09-23 — first pass, unchecked against
% the views by a human.
% Read from the local mirror only (archives/nal-1644/f0041.jpg … f0060.jpg):
% a 1000-px thumbnail of each view, then four full-width bands at 2400 px
% (normalised), and tight crops at 2–5× for struck words, signs and the
% marginal calculations. Each view was drafted as soon as it was read.
%
% Dating. The catalogue gives no dating for this volume: \dating{} is left
% empty, and this pass infers no date.
%
% What the views are. Greyscale scans, portrait, about 3750 × 5620 (4148 wide
% at view 59, where the perforated edge of a film shows at the left). One page
% per view, rectos and versos alternating: rectos at the odd views 41–59,
% versos at the even views 42–60. No view is blank; none is skipped. The
% leaves are of unequal height and are bound so that the top or the bottom of
% a neighbouring leaf projects above or below the leaf photographed, with its
% writing: views 41, 43, 44, 45, 46, 47, 49, 51, 52, 53, 54, 55, 56, 57, 58
% and 59 all show lines of another page. Those lines are transcribed only on
% the view where their own page is photographed; each view's first \note{}
% says what projects and where it belongs. Line ends and beginnings of the
% facing pages are also visible at the gutter and are ignored.
%
% What the twenty views hold. One text, in one hand — a regular, rather fast
% seventeenth-century cursive, one ink (paler on views 59–60), Latin — a
% treatise on equations in Viète's notation, set out by chapters, theorems
% and propositions, with numerical examples after each theorem. The hand
% corrects itself in the text and signals corrections in the left margin (a
% letter or a sign in the margin against an underlined word or a
% written-over letter); a few longer marginal notes, in the same hand, justify
% a step or give a variant. No title, author or date is on any leaf of the
% batch.
%
%   v. 41      Continues a text begun before the batch: the conclusion of a
%              theorem whose hypothesis is on the previous page, then Theorema
%              IX and « De quadrato quadratis adfectis sub cubo. Theorema X ».
%   v. 42      Theorem X's demonstration ends; Theorema XI; a left margin full
%              of calculations and four small figures (semicircle, right
%              triangle, lettered D C G H K L M E F).
%   v. 43–45   Theorema XII; « Deductio adfectorum aliquot cuborum ab adfectis
%              quadratis Cap. XIIII », Theoremata I–VI.
%   v. 45–46   « Ambiguitates radicum quarum potestates de homogeneis
%              adfectionum negantur demonstrata Cap. XV » (first « XIIII »,
%              struck): the two roots of B in A − A quadrato = Z plano.
%   v. 46–54   « De Syncrisi Cap. XVI »: definition of syncrisis (comparing two
%              correlated equations to find their constitution); three kinds
%              of correlated equations — ancipites, contradicentes, and a
%              third read « inversæ » with doubt; Problema I (ancipites, in
%              general, then « In Specie » for the quadratic and « Aliud » for
%              the cubic), Problema II (contradicentes), Problema III
%              (inversæ), and a closing remark on which kinds are efficacious
%              in which powers, referring to a « Caput de genesi potestatum a
%              binomia radice ».
%   v. 55–58   « Syncriticæ doctrinæ Geometrica phrasis. Cap. XVII »:
%              Propositions I–VII on series of three to seven continued
%              proportionals and their powers.
%   v. 58–60   « Æquationum ancipitum constitutiva Cap. XVIII »: « De affectis
%              sub depressiore gradu seu latere » (Props I–V), « De Adfectis
%              sub elatiore gradu seu latere recipro » (Props VI–IX), « De
%              affectis sub gradibus intermediis quæ per interpretationem
%              deprimuntur », Prop. X, whose enunciation alone stands at the
%              middle of view 60; the rest of that page is blank.
%
% Where the batch begins and ends. View 41 continues a theorem whose
% hypothesis is on the page before it (view 40, batch 2) and whose numbering
% continues the Theoremata I–VIII of a chapter that began before this batch.
% View 60 stops after the enunciation of Prop. X with the lower half of the
% page blank; whether the chapter continues on view 61 is not known from this
% batch. A leaf numbered « 29 » projects above leaf « 27 » at view 59 with
% the head of a « Prop. XV » of the same series, so the chapter does go on
% somewhere after view 60.
%
% Notation, kept as written. The unknown is A (a second root E in the
% syncrisis chapter; B D F G the continued proportionals at view 58); the
% given magnitudes are named with their dimension — « Z plano », « Z solido »,
% « B plano », « S plano plano », « Z planoplano », « Z planosolido », « Z
% solidosolido », « B coefficiens », « B parabola », « Z homogeneum » — and
% the powers in words: quadratum, cubus, quadratoquadratum, quadratocubus,
% cubocubus. « in » is multiplication; « bis », « ter » are coefficients.
% There is no sign of equality: the verb « æquatur / æquabitur » is written
% out, often after a brace. A long division is written numerator over a
% line over divisor, and is set as an array with \hline. Numerical examples
% are in the cossic abbreviations N (number, root), Q (square), C (cube), QQ,
% QC, CC, with the coefficient before: « 12N − 1Q æquatur 20 ». A sign like
% an « L » before a number marks a root, « L 20 », « LC 40 », « LQQ 80 »,
% « LQC 160 » (views 42, 58, 59); it is set $L$. The period's double stroke
% « = » appears where the argument asks for a minus (views 51, 52, 53) and is
% set « = » and noted each time; the writer's marginal signs against those
% lines suggest he corrected them. Abbreviations are kept (« quad. », « æquat. »,
% « Aq », « Eqq », « q; » for -que). The long s is set s.
%
% Numbers on the leaves. An ink number stands at the top right of each recto,
% one per leaf: 18 (v. 41), 19 (43), 20 (45), 21 (47), 22 (49), 23 (51), 24
% (53), 25 (55), 26 (57), 27 (59); versos carry none. A « 29 » is read on the
% projecting head of a later leaf at view 59; no leaf « 28 » is photographed
% in the batch, and the batch cannot say where it is. The series is one number
% per leaf, but it is in ink, not a pale pencil, and the writer's own marginal
% note at view 43 sends the reader « ad sinistram paginam fol. 18. », naming
% leaf 18 by this very number: the series is therefore taken to be the
% writer's (or a contemporary's) foliation, not the library's, and no \folio{}
% is written. Batch 2's drafts read « 9 », « 10 », « 11 » on views 22–26 in
% the same position, which fits one series; the coordinator will reconcile.
% Other numbers in the batch are content: chapter numbers (XIIII, XV, XVI,
% XVII, XVIII), theorem and proposition numbers, the ranks I–VI written above
% the terms of the numerical series, and the numbers of the examples. A « 1 »
% in a circle at view 43 is the end of « Fit 1N 21 » of view 42, seen at the
% gutter.

% Coordinator's reconciliation (2026-09-23), after all nine batches:
%   One ink series of leaf numbers runs top right of the rectos through the
%   whole volume, with no gap: 1–8 and « 5 bis » (v. 5–20), 9–17 (v. 22–39),
%   18–27 (v. 41–59), 28–37 (v. 61–79), 38–47 and « 46 bis » (v. 81–100),
%   « 47 bis » and 48–56 (v. 102–120), 57–66 (v. 122–140), 67–78
%   (v. 142–160), 79 (v. 162). It crosses the three pieces of the volume, and
%   the writer cites it himself (« fol. 18. », v. 43; « pag. 48 », v. 100), so
%   it is the writer's or copyist's count of leaves, not the library's. Being
%   ink, it gets no \folio{}, in any batch.
%   The pieces: « Francisci Vietæ De Recognitione Æquationum tractatus »
%   (heading on v. 5), Cap. 1–XXI, v. 5–68; « Francisci Vietæ. De æquationum
%   Emendatione tractatus secundus » (heading on v. 69), Cap. I–XIIII, ending
%   « Finis » on v. 141; then astronomy without a name, v. 142–163 (lunar
%   prostaphaereses, Tycho's lunar hypothesis restated, « Harmonici
%   Ptolemaici Constructio Caput XII », Mercury). Two doubtful notes may mark
%   the treatises as a copy: « transcripta ex exemplari » (v. 100) and
%   « deerant hæc in exemplari » (v. 114). The name on the leaves is in the
%   headings; who wrote the hands is not settled.

\documentclass[11pt,a4paper]{article}
% The preamble shared by every transcription of the mathematicians' manuscripts in Gallica.
%
% It rests on one idea: **uncertainty compounds, it does not resolve.** A
% transcription that smooths over a doubtful word has destroyed the only thing
% it was made for — a reader must be able to tell, at every point, what was on
% the page and what was supplied. The apparatus macros below are therefore the
% critical apparatus, not decoration.
%
% The same commands are recognised by `scripts/render.mjs`, which produces the
% site's reading view, and by `scripts/tei.mjs`. A macro added here must be
% added there too, or rendering fails loudly — which is the wanted behaviour.
%
% Comments here are English, like the rest of the repository. The documents
% this preamble sets are French, and that is deliberate: see the skills.

\ProvidesPackage{germain}[2026/09/15 Transcriptions of mathematicians' manuscripts in Gallica]

\usepackage{amsmath,amssymb,amsthm}
\usepackage{mathrsfs}
% Kept from the parent preamble so the renderer's diagram support stays
% available; these manuscripts draw no commutative diagrams, and nothing here uses it.
\usepackage{tikz-cd}
\usepackage[english,french]{babel}
\usepackage{fontspec}

% --- Greek ------------------------------------------------------------------
% The text face stays fontspec's default, Latin Modern, which has no Greek:
% XeTeX drops every glyph it lacks with a line in the log and nothing on the
% page, and the Greek words the manuscripts quote vanished from the PDFs. So
% the Greek and Greek Extended blocks get a XeTeX character class of their own,
% and entering or leaving a run of it switches the family to CMU Serif — the
% same Computer Modern design, with polytonic Greek — and back. Only the
% family changes: a Greek word in \textit or \textbf keeps its shape and series.
% The transitions to and from every other class carry the tokens that class
% has towards an ordinary letter, so French spacing before « ; » or inside
% « » still holds after a Greek word. No grouping, so no brace or \endgroup
% can come out unbalanced; maths is left alone, since XeTeX inserts nothing
% there. Kept identical in `exercices.sty`.
\makeatletter
\newfontfamily\germainGreekfont{cmun}[
  Extension=.otf, NFSSFamily=germaingreek,
  UprightFont=*rm, ItalicFont=*ti, SlantedFont=*sl,
  BoldFont=*bx, BoldItalicFont=*bi, BoldSlantedFont=*bl]
\newXeTeXintercharclass\germain@greek
\def\germain@greekrange#1#2{%
  \@tempcnta=#1\relax
  \loop
    \XeTeXcharclass\@tempcnta=\germain@greek
    \ifnum\@tempcnta<#2\advance\@tempcnta\@ne\repeat}
\germain@greekrange{"0370}{"03FF}
\germain@greekrange{"1F00}{"1FFF}
\def\germain@greekprev{\rmdefault}
\def\germain@greekin{%
  \edef\germain@greekprev{\f@family}\fontfamily{germaingreek}\selectfont}
\def\germain@greekout{\fontfamily{\germain@greekprev}\selectfont}
\def\germain@greekjoin{%
  \XeTeXinterchartokenstate=\@ne
  \@tempcnta=\z@
  \loop
    \ifnum\@tempcnta=\germain@greek\else
      \XeTeXinterchartoks\germain@greek\@tempcnta=\expandafter{%
        \expandafter\germain@greekout\the\XeTeXinterchartoks\z@\@tempcnta}%
      \XeTeXinterchartoks\@tempcnta\germain@greek=\expandafter{%
        \the\XeTeXinterchartoks\@tempcnta\z@\germain@greekin}%
    \fi
    \ifnum\@tempcnta<4095\advance\@tempcnta\@ne\repeat}
% babel-french sets its own classes, and saves and restores the interchar
% state, each time a language is selected: join again after it does.
\addto\extrasfrench{\germain@greekjoin}
\addto\extrasenglish{\germain@greekjoin}
\AtBeginDocument{\germain@greekjoin}
\makeatother

\usepackage{geometry}
\geometry{a4paper,margin=25mm,marginparwidth=28mm}
\usepackage{marginnote}
\usepackage{xcolor}
\usepackage[normalem]{ulem}
\setlength{\emergencystretch}{3em}
\usepackage{hyperref}
\hypersetup{colorlinks=true,linkcolor=black,urlcolor=black,pdfborder={0 0 0}}

% --- Batch metadata ---------------------------------------------------------
% Taken as they stand by the rendering script, which shows them at the head of
% the reading view. They fix what the file claims to cover. The macro names are
% the parent project's, so that the scripts stay shared; \folder here means the
% volume's slug — `fr-9115` — and \pages counts Gallica views.

\newcommand{\folder}[1]{\gdef\grothFolder{#1}}
\newcommand{\batch}[1]{\gdef\grothBatch{#1}}
\newcommand{\pages}[2]{\gdef\grothFirst{#1}\gdef\grothLast{#2}}
\newcommand{\foldertitle}[1]{\gdef\grothFoldertitle{#1}}

% The holder's own shelfmark, printed as they write it: « Français 9115 ».
\newcommand{\shelfmark}[1]{\gdef\grothShelfmark{#1}}

% The dating, copied verbatim from the holder's catalogue — never inferred
% here. « XIXe siècle » is the BnF's; a pass that dates a leaf from its content
% says so in a \note{}, as its own inference.
\newcommand{\dating}[1]{\gdef\grothDating{#1}}

% The status notice, printed in the title block and repeated as a diagonal
% watermark on every page of the PDF. The manuscripts are in the public
% domain; what this line declares is not a right but a quality — an
% unverified first pass by a machine. The skills write the line; a file
% without it gets no watermark, so leaving it out is visible in the output.
\newcommand{\watermark}[1]{\gdef\grothWatermark{#1}}

\gdef\grothFolder{?}\gdef\grothBatch{}\gdef\grothFirst{?}\gdef\grothLast{?}%
\gdef\grothFoldertitle{}\gdef\grothShelfmark{}\gdef\grothDating{}\gdef\grothWatermark{}

% --- The résumé -------------------------------------------------------------
\newenvironment{resume}
  {\small\setlength{\parskip}{0.45em}}
  {\par\medskip\hrule\medskip}

% --- Keywords ---------------------------------------------------------------
% In English, closing the modernised reading's résumé: the modern vocabulary
% under which to search for what the leaves construct. The manifest extracts
% them to tag the volume — this line is the single source of the tags.
\newcommand{\keywords}[1]{\par\smallskip\noindent
  {\footnotesize\textsc{Keywords} --- \textit{#1}}\par}

% --- The critical apparatus -------------------------------------------------

% The Gallica view number, in the margin. This is the anchor that lets a
% disagreement over a formula be settled by looking at *one* image rather than
% a volume of 750. The site uses it to turn the facsimile as the transcription
% is scrolled. A view is one scanned image: usually one side of a leaf.
\newcommand{\page}[1]{%
  \marginnote{\footnotesize\color{gray!70}\textsf{v.\,#1}}%
  \label{page:#1}\ignorespaces}

% The BnF's pencilled foliation, where the leaf carries it and a pass can read
% it: « 198r ». This is what the literature cites — « ff. 198r–208v » — and
% the one line that turns a citation into a view. Recorded, never inferred:
% a folio a pass cannot read is not written.
\newcommand{\folio}[1]{%
  \marginnote{\footnotesize\color{gray!50}\textsf{f.\,#1}}\ignorespaces}

% The modernised reading's coarser counterpart to \page{}: which views a
% section is drawn from.
\newcommand{\pagerange}[2]{%
  \marginnote{\footnotesize\color{gray!70}\textsf{v.\,#1--#2}}%
  \ignorespaces}

% Illegible: never guessed.
\newcommand{\ill}{\textcolor{red!60!black}{[\,\ldots\,]}}

% A doubtful reading: the word is given, and flagged as uncertain.
\newcommand{\uncertain}[1]{\uline{#1}}

% An editorial addition — a missing letter, an implied word.
\newcommand{\add}[1]{\textcolor{blue!50!black}{[#1]}}

% Struck out by the writer. What was crossed out is part of the document.
\newcommand{\struck}[1]{\sout{#1}}

% The transcriber's note, on the state of the page or the mathematics.
\newcommand{\note}[1]{\par\smallskip{\small\color{gray!60!black}\textit{[#1]}}\par\smallskip}

% A marginal note of hers, distinct from the body of the text.
\newcommand{\marginal}[1]{\par\smallskip\noindent\hspace{1em}%
  {\small\color{gray!60!black}#1}\par\smallskip}

% --- Title ------------------------------------------------------------------
% The title block is French, like the documents themselves.

\AddToHook{shipout/background}{%
  \ifx\grothWatermark\empty\else
    \begin{tikzpicture}[remember picture, overlay]
      \node[rotate=52, scale=3.4, align=center, text=gray!18,
            font=\sffamily\bfseries]
        at (current page.center) {\grothWatermark};
    \end{tikzpicture}%
  \fi}

\AtBeginDocument{%
  \begin{center}
    {\large\bfseries Manuscrits de mathématiciens (Gallica) ---
      \ifx\grothShelfmark\empty\grothFolder\else\grothShelfmark\fi}\\[2pt]
    {\itshape\grothFoldertitle}\\[4pt]
    {\small vues \grothFirst--\grothLast
      \ifx\grothBatch\empty\else\ (lot \grothBatch)\fi}\\[2pt]
    \ifx\grothDating\empty\else
      {\small\color{gray!60!black}Datation du catalogue : \grothDating}\\[2pt]
    \fi
    \ifx\grothWatermark\empty\else
      {\small\bfseries\color{red!45!gray}\grothWatermark}\\[2pt]
    \fi
    {\footnotesize\color{gray!60!black}Transcription établie d'après les images IIIF de Gallica.
     Source gallica.bnf.fr / Bibliothèque nationale de France.\\
     Les lectures incertaines sont soulignées, les illisibles notées [\,\ldots\,],
     les ajouts entre crochets ; \textsf{v.} numérote les vues de Gallica,
     \textsf{f.} la foliotation au crayon de la BnF.}
  \end{center}
  \vspace{1em}}

\endinput


\folder{nal-1644}
\shelfmark{NAL 1644}
\batch{3}
\pages{41}{60}
\dating{}
\watermark{Lecture automatique\\non vérifiée}
\foldertitle{Viète (François). Papiers divers. NAL 1644}

\begin{document}

\page{41}
\note{Recto. En haut à droite, à l'encre, le chiffre « 18 » (voir l'en-tête pour la série). Le feuillet est plus court que le suivant : sous son bord inférieur, irrégulier, dépasse le bas du feuillet chiffré « 19 », avec « $-$ A cubo », « ut est enunciatum », l'exemple « Si 1Q + 8N æquet. 20 / 84N $-$ 1C æquabit 160 » et, en marge, « Fit 1N 2. » ; ces lignes appartiennent à la page de la vue 43 et y sont transcrites. À gauche, au pli, restent visibles des fins de lignes de la page précédente (« …remate », « …plano », « ergo »), qui ne sont pas de cette page. La page commence au milieu d'un théorème dont l'hypothèse est à la page précédente : elle en donne la conclusion.}

\[ \left.\begin{array}{l} A \text{ quad.quad.} \\ \left.\begin{array}{l} + D \text{ plano bis} \\ - B \text{ quad.} \end{array}\right\} \text{in } A \text{ quadrat.} \\ - B \text{ in } S \text{ plano bis in } A \end{array}\right\} \text{æquabitur} \left\{\begin{array}{l} S \text{ plano plano} \\ - D \text{ plano plano.} \end{array}\right. \]
\note{Devant « B in S plano bis in A », entre le signe $-$ et le B, un signe en forme de dièse (« \(\#\) ») : une lettre biffée de deux traits croisés, sans doute ; elle n'est pas lue.}

Nec dissimilis est ab ea quæ in antecedente theoremate exposita est demonstratio.

\marginal{Fit 1N. 10.}

Si 1Q $-$ 8N. æquetur 20 differentiæ inter 40 et 60.

1QQ + 16Q $-$ 960N. æquabitur 2000.
\note{« Fit 1N. 10. » est écrit dans la marge de gauche, en regard de l'exemple : la racine de l'exemple numérique. Les puissances des nombres sont notées N (le nombre, la racine), Q (le carré), C (le cube), QQ (le carré-carré), précédées de leur coefficient.}

\subsection*{Theorema IX.}

\[ \text{Si } \left.\begin{array}{l} B \text{ in } A \\ - A \text{ quadrato} \end{array}\right\} \text{æquetur} \left\{\begin{array}{l} S \text{ plano} \\ - D \text{ plano.} \end{array}\right. \]

\[ \left.\begin{array}{l} B \text{ in } D \text{ plano bis in } A \\ \left.\begin{array}{l} + B \text{ quadrato} \\ - S \text{ plano bis} \end{array}\right\} \text{in } A \text{ quadrato} \\ - A \text{ quad.quad.} \end{array}\right\} \text{æquabitur} \left\{\begin{array}{l} S \text{ plano plano} \\ - D \text{ plano plano} \end{array}\right. \]
\note{La lettre de « $-$ S plano bis » est écrite sur une autre, barrée d'une croix ; en marge de gauche, en regard, « S/ » : la correction est signalée en marge. La lettre première n'est pas lue (D ?).}

Nec dissimilis est ab ea quæ in \struck{antecedente} septimo huius theorematis capite exposita est demonstratio.
\note{« antecedente » est biffé d'un trait et « septimo » écrit à la suite, sur la même ligne. « theorematis capite » : lecture des lettres.}

\marginal{Fit 1N. 2.}

Si 12N. $-$ 1Q æquetur 20 differentiæ inter 40 et 60.

960N. + 24Q $-$ 1QQ æquabitur 2000.
\note{Avec la racine 2 donnée en marge, les deux égalités numériques sont justes ($24 - 4 = 20$ ; $1920 + 96 - 16 = 2000$) ; elles supposent $D = 40$ dans « B in D plano bis » et $S = 60$ dans « S plano bis », tandis que le second membre, $2000 = 60^2 - 40^2$, est écrit « S plano plano $-$ D plano plano ». Le transcripteur ne tranche pas.}

\subsection*{De quadrato quadratis adfectis sub cubo. Theorema X.}

\[ \text{Si } \left.\begin{array}{l} A \text{ quadrato} \\ + B \text{ in } A \end{array}\right\} \text{æquetur } Z \text{ plano} \]

\[ \left.\begin{array}{l} \begin{array}{l} A \text{ quad.quad.} \\ \left.\begin{array}{l} + B \text{ in } Z \text{ plano bis} \\ + B \text{ cubo} \end{array}\right\} \text{in } A \text{ cubum} \end{array} \\ \hline \begin{array}{l} Z \text{ plano} \\ + B \text{ quadrato} \end{array} \end{array}\right\} \text{æquabitur } \begin{array}{c} Z \text{ plano plano} \\ \hline \begin{array}{l} Z \text{ plano} \\ + B \text{ quadrato.} \end{array} \end{array} \]
\note{Les deux membres sont écrits chacun au-dessus d'un trait, sous lequel est écrit « Z plano + B quadrato » : une division. « in A cubum » est la lecture des lettres, nette sur le feuillet ; ni l'énoncé ni son rapport avec la démonstration qui suit (laquelle conclut sur la valeur de A) ne sont corrigés ici.}

quoniam enim $A$ quadratum plus $B$ in $A$ æquatur $Z$ plano. igitur per antithesim $A$ quadratum æquabitur $Z$ plano $-$ $B$ in $A$. omnia ducantur in $A$. ergo $A$ cubus æquat. $Z$ plano in $A$ $-$ $B$ in $A$ quadratum. seu aliter ex designato $A$ quadrati valore

\[ A \text{ cubus æquat. } \begin{array}{l} Z \text{ plano in } A \\ - B \text{ in } Z \text{ plano} \\ + B \text{ quadrato in } A \end{array} \]

\marginal{quoniam $A$ quad. \ill{} $Z$ p. \ill{} et $A$ c. \(\}\) $Z$ p\ill{} $-$ $B$ in $A$ q. \ill{} in affectione $-$ $B$ in $A$ q. pro $A$ q. sumatur $Z$ p. $-$ $B$ in $A$ fiet $A$ cub. \(\{\) $Z$ pla. in $A$ $-$ $B$ in $Z$ pl. $+$ $B$ quad. in $A$.}
\note{Note de la marge de gauche, de la même main et de la même encre, en regard des lignes précédentes ; elle refait le passage de « $A$ cubus æquat. $Z$ plano in $A$ $-$ $B$ in $A$ quadratum » à la forme ci-dessus en remplaçant $A$ quadratum par sa valeur. Une grosse tache d'encre allongée, en diagonale, couvre la fin de ses deux premières lignes ; ce qu'elle cache n'est pas lu. Le mot après « $-$ B in A q. », à la fin de la deuxième ligne, est \ill{} (« posito » ? les lettres ne le portent pas). Le début de la troisième ligne, contre le bord, est lu « in » avec doute.}

Et utraq; æquationis huius parte divisa per $Z$ plano $+$ $B$ quad.

\[ \left.\begin{array}{c} \begin{array}{l} A \text{ cubus} \\ + Z \text{ plano in } B \end{array} \\ \hline \begin{array}{l} Z \text{ plano} \\ + B \text{ quadrato} \end{array} \end{array}\right\} \text{æquatur } A. \]
\note{Fin du texte du recto, au bord inférieur du feuillet ; la démonstration continue en tête du verso (vue 42). Ce qui dépasse sous ce bord appartient au feuillet chiffré « 19 » (vue 43). Dans l'énoncé ci-dessus, le second membre est écrit « Z plano plano » ; la même égalité, à la fin de la démonstration (vue 42), porte « Z plano plano plano ».}

\page{42}
\note{Verso du feuillet chiffré « 18 » ; aucun chiffre en tête. En haut, au-dessus de la première ligne, un astérisque à huit branches, à l'encre, sans renvoi repéré sur la page. À droite de la vue dépassent les bords des feuillets suivants, avec des débuts de lignes qui ne sont pas de cette page. Le texte continue sans rupture la démonstration du théorème X laissée au bas de la vue 41.}

quod igitur diximus $A$ quadratum æquari $Z$ plano $-$ $B$ in $A$. id est ex designato $A$ valore ita exprimitur

\[ A \text{ quadratum æquatur } \begin{array}{c} \begin{array}{l} Z \text{ plano plano} \\ - B \text{ in } A \text{ cubum} \end{array} \\ \hline \begin{array}{l} Z \text{ plano} \\ + B \text{ quadrato} \end{array} \end{array} \]
\note{Le numérateur est écrit sur deux lignes, « Z plano plano » puis « $-$ B in A cubum », et le trait de division est tiré sous la seconde seulement, un peu plus court que la première.}

Rursus iterum venio ad primam propositam æquationem

\[ \left.\begin{array}{l} A \text{ quadratum} \\ + B \text{ in } A \end{array}\right\} \text{æquari } Z \text{ plano} \]

utraq; pars ducitur quadraticè. fiet

\[ \left.\begin{array}{l} A \text{ quadrato quadratum} \\ + B \text{ in } A \text{ cubum bis} \\ + B \text{ quadrato in } A \text{ quadrat.} \end{array}\right\} \text{æquabitur } Z \text{ plano plano} \]

Iam affectio $B$ quadrati in $A$ quadratum interpretationem accipito ex postremo designato $A$ quadrati valore. Ipsa erit

\[ \begin{array}{c} \begin{array}{l} B \text{ quad. in } Z \text{ plano plano} \\ - B \text{ cubo in } A \text{ cubum} \end{array} \\ \hline \begin{array}{l} Z \text{ plano} \\ + B \text{ quadrato} \end{array} \end{array} \]

qua interpretatione adscita et omnibus rite ordinatis

\[ \left.\begin{array}{l} A \text{ quadrato quadratum} \\ \begin{array}{c} \left.\begin{array}{l} + B \text{ in } Z \text{ plano bis} \\ + B \text{ cubo} \end{array}\right\} \text{in } A \text{ cubum} \\ \hline \begin{array}{l} Z \text{ plano} \\ + B \text{ quadrato} \end{array} \end{array} \end{array}\right\} \text{æquabitur } \begin{array}{c} Z \text{ plano plano plano} \\ \hline \begin{array}{l} Z \text{ plano} \\ + B \text{ quadrato} \end{array} \end{array} \]

ut est ordinatum.
\note{Ici le second membre est « Z plano plano plano », là où l'énoncé du théorème X, à la vue 41, porte « Z plano plano » : les deux sont gardés tels qu'ils sont écrits ; l'exemple qui suit ($147^3 / 343 = 9261$) s'accorde avec le premier. « ordinatum » : lecture des lettres, là où la vue 41 écrit, à la même place d'une démonstration, « enunciatum ».}

Si 1Q + 14N æquat. 147

1QQ + 20C æquabit 9261.

\marginal{Fit 1N. 7.}
\note{« Fit 1N. 7. » est écrit à droite de l'exemple et entouré d'un trait.}

\subsection*{Theorema XI.}

\[ \text{Si } \left.\begin{array}{l} A \text{ quadratum} \\ - B \text{ in } A \end{array}\right\} \text{æquetur } Z \text{ plano} \]

\[ \left.\begin{array}{c} \begin{array}{l} A \text{ quadrat. quadrat.} \\ \left.\begin{array}{l} - B \text{ in } Z \text{ plano bis} \\ - B \text{ cubo} \end{array}\right\} \text{in } A \text{ cubum} \end{array} \\ \hline \begin{array}{l} Z \text{ plano} \\ + B \text{ quadrato} \end{array} \end{array}\right\} \text{æquabitur } \begin{array}{c} Z \text{ plano plano plano} \\ \hline \begin{array}{l} Z \text{ plano} \\ + B \text{ quadrato.} \end{array} \end{array} \]
\note{Comme au théorème X, le trait de division est tiré sous tout le premier membre, « A quadrat. quadrat. » compris.}

Nec dissimilis est ab ea quæ in antecedente theoremate exposita est demonstratio

Si 1Q $-$ 14N. æquetur 147

1QQ $-$ 20C æquabitur 9261.

\marginal{Fit 1N 21}
\note{« Fit 1N 21 » est écrit à droite et entouré d'un trait. La suite de la page, sous l'exemple, est blanche ; au pied, le bord d'un feuillet plus grand dépasse, sans écriture.}

\marginal{$D G q$. $-$ $E C q$. $|$ $D C q$. $|$ $C G q$ $|$ $E C$ \quad $E C$ cub. \quad $C G$ in $D C q$. \quad æqualit \quad $D G q$ in $E C$ \quad 6 in 6 [8] \quad 104 \quad $L$ 52. in 2. \quad $\frac{52}{36}$ $|$ 1 $\frac{16}{36}$ \quad $L \frac{81}{169}$ $-$ 52 \quad 13 \quad 9 \quad 52. 36. $L$ 52 \quad 52 \quad 16 \quad $L$ 52 \quad 13 \quad $C G$ $64$ $-$ $C H$. \quad $Z$ plano plano $+$ $B$ quad. $|$ $Z$ plano bis $+$ $B$ q. $|$ $B$ \quad $D G$ quad. $|$ $D G$ quad. $+$ $C G$ q. $|$ $C G$ $|$ $C G$ $C H$ \quad $Z$ plano $+$ $B$ quad. \quad $Z$ plano $Z$ plano plano. \quad $D G$ q. \quad $D C$ q. \quad $D q$. $q$. \quad $D K q$. \quad 52 \quad 16 \quad 256 \quad $D g$. $|$ $C G$ q. $D G$ q. $|$ $D G$ $|$ $H G$ $D G$ \quad 52 \quad 108 \quad $D G$ \quad $D C$ \quad $C G$ \quad $C m$ \quad 6 \quad 12 $\frac{15}{26}$}
\note{La marge de gauche, sur toute la hauteur de la page, porte de la même encre des calculs et quatre figures, qui traduisent en géométrie les analogies du texte ; ils sont donnés ci-dessus dans l'ordre de haut en bas, sans prétendre en rendre la disposition, qui est faite de colonnes séparées de traits verticaux, de traits de division et d'accolades. Figures : (1) en haut, un demi-cercle sur le diamètre $E\,F$, les points $E$, $C$, $G$, $F$ sur le diamètre ; de $D$, au sommet, la perpendiculaire $D\,C$ et la droite $D\,G$ prolongée au-dessous du diamètre jusqu'en $L$ ; sur $D\,C$ le point $K$, d'où une parallèle au diamètre coupe $D\,G$ en $H$ ; $C\,H$ joint ; au-dessus du diamètre, près de $F$, la lettre $M$, d'où une perpendiculaire descend vers $L$. (2) Plus bas, un demi-cercle de même construction, sur $E$, $C$, $A$, $G$, $H$, chiffré : « 2 », « 3 », « 3 », « 2 » le long du diamètre, « 4 » et « 5 » sur les droites, « 169 », « 16 », « 52 », « 89 » (lecture douteuse) et « $D\,C\,q = 16$ » près de l'arc. (3) Au milieu, un petit triangle rectangle $D\,C\,G$, avec $K$ sur $D\,C$ et $H$ sur $D\,G$. (4) En bas, le même triangle chiffré « 512 », « 640 », « 800 », « 480 », « 360 », « 600 », « 216 », « 360 », avec un nombre surchargé au sommet. Plusieurs lectures de ces calculs sont incertaines : « [8] » est un 8 surmonté d'un trait, au-dessus de « 6 in 6 » ; « $L$ » est un signe en forme de L, qui paraît noter une racine (« $L$ 52 ») ; « $\frac{81}{169}$ » peut se lire « 51 » ; « $D q$. $q$. » est un « q » écrit sous un autre.}

\page{43}
\note{Recto. En haut à droite du feuillet, à l'encre, « 19 ». Au-dessus de son bord supérieur dépasse le haut d'un feuillet plus grand, chiffré « 21 », avec trois lignes et demie de texte (« Sive autem ancipitum, sive contradicentium, sive inversarum, æquationum constitutio… ») : c'est la tête de la page de la vue 47, où elles sont transcrites. À gauche de la vue, au pli, des fins de lignes de la page précédente. Le texte de ce recto ne reprend pas celui du verso de la vue 42, qui finissait sur l'exemple du théorème XI : il continue la suite des théorèmes avec le théorème XII, puis ouvre un nouveau chapitre.}

\subsection*{Theorema XII.}

\[ \text{Si } \left.\begin{array}{l} B \text{ in } A \\ - A \text{ quadrato} \end{array}\right\} \text{æquetur } Z \text{ plano.} \]

\[ \left.\begin{array}{c} \begin{array}{l} \left.\begin{array}{l} B \text{ cubus} \\ - B \text{ in } Z \text{ plano bis} \end{array}\right\} \text{in } A \text{ cub.} \\ - A \text{ quad. quadrato.} \end{array} \\ \hline \begin{array}{l} B \text{ quadrato} \\ - Z \text{ plano} \end{array} \end{array}\right\} \text{æquabitur } \begin{array}{c} Z \text{ plano plano plano} \\ \hline \begin{array}{l} B \text{ quadrato} \\ - Z \text{ plano} \end{array} \end{array} \]
\note{La formule est très corrigée. À gauche, « $-$ A quad. quadrato. » est écrit en interligne, au-dessus d'une ligne biffée qui se lit « $-$ \ill{} quadrato in A cubum » ; une accolade courbe enveloppe les lignes du premier membre, et un trait courbe rattache « in A cub. » aux deux premières. Sous chacun des deux traits de division, « B quadrato » et « $-$ Z plano » sont surchargés (le $-$ de droite sur un $+$), et une troisième ligne, biffée, se lit « \ill{} quadrato ». Dans la marge de droite, encadrée d'un trait, la même égalité récrite au net en abrégé : « Bc. / $-$ B in Z p. bis \(\}\) in A c. » sur un trait, « Bq / $-$ Z p. » dessous, puis, sous un second trait, « $-$ A qq. ». Le transcripteur donne l'état final lisible ; il n'est pas sûr que « A quad. quadrato » soit, dans l'intention de l'écrivain, hors de la division, comme le porte la note de marge.}

Nec dissimilis est ab ea quæ in decimo huius capitis theoremate exposita est demonstratio

Si 21N $-$ 1Q æquetur 98

15C $-$ 1QQ æquabitur \uncertain{2}744
\note{Le premier chiffre de « 2744 » est surchargé : un 7 et un 2 l'un sur l'autre. Avec la racine 7 donnée en marge, $15 \cdot 343 - 2401 = 2744$. Dans la marge de gauche, en regard de cette ligne, une tache d'encre et un trait vertical.}

\marginal{Fit 1N 7}

\marginal{$\leftarrow$ ad sinistram paginam fol. 18. quoniam $A$ c. $+$ $Z$ p in $B$ sur $Z$ p. $+$ $B$ q. \(\}\) $A$ ergo $B$ in $A$ erit $B$ in $A$ c. $+$ $B$ q in $Z$ p sur $Z$ p. $+$ $B$ q. quo a $Z$ pla. ablato remanebit $Z$ pp. $-$ $B$ in $A$ c. sur $Z$ p. $+$ $B$ q. \quad \ill{} $B$ in $Z$ pp. $-$ $B$ in $A$ c. sur $Z$ p. $+$ $B$ q.}
\note{Note de la marge de gauche, de la même main, en deux morceaux enfermés chacun dans un trait courbe ; une flèche en tête pointe vers la page de gauche. Elle justifie la première égalité de la vue 42 (« A quadratum æquatur Z plano plano $-$ B in A cubum » sur « Z plano $+$ B quadrato ») à partir de la valeur de A trouvée à la vue 41. « sur » rend ici un trait de division : chaque fraction est écrite, numérateur et diviseur, l'un sous l'autre. Le renvoi « fol. 18. » nomme le feuillet par le chiffre qui se lit en tête de la vue 41. Le second morceau commence par deux ou trois mots \ill{} (« vt ha… » ?).}

\subsection*{Deductio adfectorum aliquot cuborum ab adfectis quadratis Cap. XIIII Theorema I}
\note{Titre centré, sur trois lignes. Le numéro du chapitre est écrit « Xiiii ».}

\[ \text{Si } \left.\begin{array}{l} A \text{ quadratum} \\ + B \text{ in } A. \end{array}\right\} \text{æquetur } Z \text{ plano} \]

\[ \left.\begin{array}{l} \left.\begin{array}{l} B \text{ quadratum} \\ + Z \text{ plano} \end{array}\right\} \text{in } A \\ - A \text{ cubo} \end{array}\right\} \text{æquabitur } B \text{ in } Z \text{ plano.} \]

quoniam enim
\[ \left.\begin{array}{l} A \text{ quadratum} \\ + B \text{ in } A \end{array}\right\} \text{æquatur } Z \text{ plano} \]
omnia ducantur in $A$. igitur
\[ \left.\begin{array}{l} A \text{ cubus} \\ + B \text{ in } A \text{ quadratum} \end{array}\right\} \text{æquabitur } Z \text{ plano in } A. \]
Sed et ex iis quæ proposita sunt
\[ A \text{ quadrat. æquat. } \left\{\begin{array}{l} Z \text{ plano} \\ - B \text{ in } A. \end{array}\right. \]
quare adscita ea interpretatione in æquatione cubica ad exprimendum valorem solidi $B$ in $A$ quadratum.
\[ \left.\begin{array}{l} A \text{ cubus} \\ + B \text{ in } Z \text{ plano} \\ - B \text{ quad. in } A \end{array}\right\} \text{æquabitur } Z \text{ plano in } A. \]
Et omnibus rite ordinatis
\[ \left.\begin{array}{l} \left.\begin{array}{l} Z \text{ plano} \\ + B \text{ quadrato} \end{array}\right\} \text{in } A \\ - A \text{ cubo} \end{array}\right\} \text{æquabitur } B \text{ in } Z \text{ plano.} \]
ut est enunciatum.

\marginal{Fit 1N 2.}

Si 1Q + 8N. æquet. 20

84N. $-$ 1C æquabit 160.
\note{Fin de la page, au pied du feuillet. C'est ce bas de page qui dépasse sous le feuillet « 18 » à la vue 41. À l'extrême gauche de la vue, hors du feuillet, « …ita » et un « 1) » dans un trait courbe : les fins de « exposita » et de « Fit 1N 21 » entouré, de la page de la vue 42.}

\page{44}
\note{Verso du feuillet « 19 » ; aucun chiffre en tête. Le feuillet est plus court que le feuillet « 18 » : au-dessus de son bord supérieur dépasse le haut du verso de la vue 42 (l'astérisque, le haut de la première figure de la marge, « quod igitur diximus A quadrat. æquari Z plano $-$ B in A. id est ex designato A valore ita exprimitur »), et, à gauche, la marge de calculs de la même page (« Dg », « 52 »…) ; ces lignes sont transcrites à la vue 42. À droite, au pli, les débuts de lignes d'une autre page (« fit 1… »). La page suit le recto (vue 43) : théorèmes II à IIII du chapitre ouvert au recto.}

\subsection*{Theorema II}

\[ \text{Si } \left.\begin{array}{l} A \text{ quadrat.} \\ - B \text{ in } A \end{array}\right\} \text{æquetur } Z \text{ plano} \]

\[ \left.\begin{array}{l} A \text{ cubus} \\ \left.\begin{array}{l} - B \text{ quadrato} \\ - Z \text{ plano} \end{array}\right\} \text{in } A \end{array}\right\} \text{æquabitur } B \text{ in } Z \text{ plano} \]

Nec dissimilis est ab ea quæ in antecedente theoremate exposita est demonstratio.

\marginal{Est 1N. 10.}

Si 1Q $-$ 8N. æquetur 20

1C $-$ 84 æquabit 160.
\note{« 1C $-$ 84 » : le signe N manque après 84 ; avec la racine 10 donnée en marge, $1000 - 840 = 160$.}

\subsection*{Theorema III}

\[ \text{Si } \left.\begin{array}{l} B \text{ in } A \\ - A \text{ quadrato} \end{array}\right\} \text{æquetur } Z \text{ plano} \]

\[ \left.\begin{array}{l} \left.\begin{array}{l} B \text{ quadratum} \\ - Z \text{ plano} \end{array}\right\} \text{in } A \\ - A \text{ cubo} \end{array}\right\} \text{æquabitur } B \text{ in } Z \text{ plano} \]

Nec dissimilis est ab ea quæ in primo huius capitis theoremate exposita est demonstratio

\marginal{Fit 1N. 2.}

Si 12N $-$ 1Q æquet. 20

124N. $-$ 1C æquabit 240

\subsection*{Theorema IIII}

\[ \text{Si } \left.\begin{array}{l} A \text{ quadratum} \\ + B \text{ in } A \end{array}\right\} \text{æquetur } B \text{ in } D \]

\[ \left.\begin{array}{l} A \text{ cubus} \\ + B + D \text{ in } A \text{ quadrat.} \end{array}\right\} \text{æquabit. } B \text{ in } D \text{ quadratum} \]

quoniam enim $A$ quadratum $+$ $B$ in $A$ æquat. $B$ in $D$ omnia ducantur in $A$. igitur
\[ \left.\begin{array}{l} A \text{ cubus} \\ + B \text{ in } A \text{ quadrat.} \end{array}\right\} \text{æquabit. } B \text{ in } D \text{ in } A \]
Sed et ex iis quæ proponuntur
\[ \begin{array}{c} \begin{array}{l} B \text{ in } D \\ - A \text{ quadrato} \end{array} \\ \hline B \end{array} \text{ æquat. } A \]
\note{Le trait final de « proponuntur » se prolonge juste au-dessus du signe $-$, de sorte que l'ensemble a l'air d'un double trait ; le signe lui-même est simple.}
quare ea interpretatione adscita ad exprimendum valorem solidi $B$ in $D$ in $A$.
\[ \left.\begin{array}{l} A \text{ cubus} \\ + B \text{ in } A \text{ quadratum} \end{array}\right\} \text{æquabitur} \left\{\begin{array}{l} B \text{ in } D \text{ quadrat.} \\ - D \text{ in } A \text{ quadratum.} \end{array}\right. \]
\note{Le « B » de « $+$ B in A quadratum » est repassé à l'encre plus noire.}
et adhibita congrua metathesi
\[ \left.\begin{array}{l} A \text{ cubus} \\ + B + D \text{ in } A \text{ quadrat.} \end{array}\right\} \text{æquabitur } B \text{ in } D \text{ quadrat.} \]
ut est ordinatum.

\marginal{Est 1N. 4.}

Si 1Q + 16N. æquat. 80 facto ex 16 in 5.

1C + 21Q æquabit 400
\note{Fin du texte ; le bas du feuillet est blanc. Avec la racine 4 donnée en marge : $16 + 64 = 80$ et $64 + 336 = 400 = 16 \cdot 25$. À l'exemple du théorème III, avec la racine 2 : $24 - 4 = 20$ et $124 \cdot 2 - 8 = 240$ ; « 124 » est $B^2 - Z$ ($144 - 20$) et $240 = 12 \cdot 20$.}

\page{45}
\note{Recto. En haut à droite du feuillet, à l'encre, « 20 ». Au-dessus de son bord supérieur dépasse encore le haut du feuillet « 21 », avec les mêmes trois lignes et demie qu'à la vue 43 (« Sive autem ancipitum… »), transcrites à la vue 47. À gauche de la vue, au pli, des fins de lignes du verso précédent (« …exposita », « …theoremate », « …ducantur »). La page continue la série des théorèmes du chapitre ouvert à la vue 43, puis ouvre un nouveau chapitre.}

\subsection*{Theorema V.}

\[ \text{Si } \left.\begin{array}{l} A \text{ quadrat.} \\ - B \text{ in } A \end{array}\right\} \text{æquetur } B \text{ in } D \]

\[ \left.\begin{array}{l} B + D \text{ in } A \text{ quadratum} \\ - A \text{ cubo} \end{array}\right\} \text{æquabitur } B \text{ in } D \text{ quadratum.} \]

Nec dissimilis est ab ea quæ in antecedente theoremate exposita est demonstratio

\marginal{Fit 1N. 20.}

Si 1Q $-$ 16N æquet. 80 facto ex 16 in 5.

21Q $-$ 1C. æquabit 400.

\subsection*{Theorema VI}

\[ \text{Si } \left.\begin{array}{l} B \text{ in } A \\ - A \text{ quadrato} \end{array}\right\} \text{æquetur } B \text{ in } D. \]

\[ \left.\begin{array}{l} B - D \text{ in } A \text{ quadratum} \\ - A \text{ cubo} \end{array}\right\} \text{æquabitur } B \text{ in } D \text{ quadrato} \]

Nec dissimilis est quæ in quarto huius capitis theoremate exposita est demonstratio

\marginal{Fit 1N. 3.}

Si 9N $-$ 1Q æquetur 18 facto ex 9 in 2.

7Q $-$ 1C. æquabit 36.
\note{Le Q de « 7Q » est écrit sur une autre lettre ; en marge de gauche, en regard, « /Q/ » entre deux traits obliques, comme le « S/ » de la vue 41 : la correction est signalée en marge. Avec la racine 3 : $27 - 9 = 18$ et $63 - 27 = 36$.}

\subsection*{Ambiguitates radicum quarum potestates de homogeneis adfectionum \uncertain{in adæquationibus} negantur. demonstrata Cap. \struck{XIIII} XV}
\note{Titre centré, sur quatre lignes. « in adæquationibus » : le groupe initial se lit « m ad » devant « æquationibus », et le texte qui suit écrit « in æquationibus ». Le numéro du chapitre était d'abord « XIIII », biffé d'un trait, et « XV » est écrit à la suite.}

Ostendendum est cum in æquationibus potestates de homogeneis adfectionum negantur radices esse ancipites.

Proponatur differentiam inter $B$ et $A$ esse æqualem $S$. et sit $B$ maior quam $S$. aut igitur excessus est penes $B$, vel penes $A$. primo casu $B - A$ æquatur $S$. itaq; $B - S$ fit $A$.

Secundo casu $A - B$ æquat. $S$. itaq; $B + S$ fit $A$. quoniam autem primo casu $B - A$ æquat. $S$, utraq; pars æqualitatis ducatur quadraticè fiet
\[ \left.\begin{array}{l} B \text{ quadratum} \\ - B \text{ in } A \text{ bis} \\ + A \text{ quadrato} \end{array}\right\} \text{æquabitur } S \text{ quadrato.} \]
Et æqualitate ordinata
\[ \left.\begin{array}{l} B \text{ in } A \text{ bis} \\ - A \text{ quadrato} \end{array}\right\} \text{æquale fit} \left\{\begin{array}{l} B \text{ quadrato} \\ - S \text{ quadrato} \end{array}\right. \]
quoniam vero secundo casu $A - B$ æquatur $S$, utraq; pars æqualitatis ducat. quadraticè
\[ \text{fiet } \left.\begin{array}{l} B \text{ quadratum} \\ - B \text{ in } A \text{ bis} \\ + A \text{ quadrato} \end{array}\right\} \text{æquabit. } S \text{ quadrato} \]
\note{« adfectionum » (troisième ligne du texte) : la fin du mot est un simple trait. Le bas du feuillet est blanc sous ces lignes ; l'exposé continue au verso (vue 46).}

\page{46}
\note{Verso du feuillet « 20 » ; aucun chiffre en tête. Comme à la vue 44, au-dessus du bord supérieur dépasse le haut du verso de la vue 42 (l'astérisque, la première figure de la marge, les deux premières lignes et le début de la troisième), transcrit à la vue 42. La page continue sans rupture le recto (vue 45), qui finissait sur l'égalité du second cas.}

Et æqualitate ordinata
\[ \left.\begin{array}{l} B \text{ in } A \text{ bis} \\ - A \text{ quadrato} \end{array}\right\} \text{rursus æquale fit} \left\{\begin{array}{l} B \text{ quadrato} \\ - S \text{ quadrato.} \end{array}\right. \]

Utroq; igitur casu in eandem inciditur æquationis formulam: est autem radix duplex. ipsa vero formula æquationis est ut planum sub latere adficiat multa quadrati.

Sit $B$ 6. $S$ 4. $A$ 1N.

12N. $-$ 1Q æquat. 20. et fit 1N. 6 $-$ 4 vel 6 + 4.

Quam amphiboliam in omnibus similibus æquationibus locum habere satis intelligitur. ut si proponatur $\left.\begin{array}{l} D \text{ in } A \\ - A \text{ quadrato} \end{array}\right\}$ æquari $Z$ plano ipsa $D$ dicetur esse $B$ bis et $Z$ planum esse $S$ quadratum minus $B$ qu\add{adrato}
\note{« minus » est abrégé. La ligne va jusqu'au pli et « B qu… » s'y perd ; le reste du mot n'est pas visible. L'égalité ci-dessus donne pour le plan « B quadrato $-$ S quadrato » ; la phrase, telle qu'elle se lit, nomme les deux carrés dans l'ordre inverse.}

Porro per sextum theorema antecedentis capitis, apparet ab ea quadra\add{ti} ambigui æquatione deduci cubos negatos de solidi parodico sub gradu et per tertium sextum et nonum capitis noni ab eadem deduci quadrato quadrata negata de plano plano sub parodico gradu. quod ad ulterioris ordinis æquationes posse extendi satis fit manifestum.
\marginal{negatos de solidis sub parodico gradu}
\marginal{duodecimi}
\note{« per » est écrit en interligne au-dessus de « sextum », sans signe d'insertion. « quadra » est au bout de la ligne, au pli. « solidi parodico » et « sub gradu », à cheval sur deux lignes, sont soulignés ; en regard, dans la marge de gauche, de la même main, « negatos de solidis sub parodico gradu » : une nouvelle rédaction de la tournure, sans que le texte soit biffé. « noni » est souligné de même, et la marge porte en regard « duodecimi » : une correction du renvoi, sans rature du texte. Les deux leçons sont données.}

\subsection*{De Syncrisi Cap. XVI}

Syncrisis est duarum æquationum correlatarum mutua inter se ad deprehendendam earum constitutionem collatio.

Duæ autem æquationes correlatæ intelliguntur cum ambæ similes sunt et præterea iisdem datis magnitudinibus constant sive adfectionum parabolis sive adfectionis homogeneis, radices tamen ideo diversæ sunt quoniam vel ipsæ æquationis formulæ de duabus pluribusve radicibus ex sui constitutione sunt explicabiles vel in iis diversa est adfectionis qualitas seu nota.

At de simplicioribus quidem correlatis sufficiet doctrina id est una tantum adfectione obnoxiis, ut pote quarum vi, de iis quæ \uncertain{passionibus} abundat, \uncertain{plasticas} potius ratiocinabitur \ill{}.
\note{« passionibus » : le milieu du mot est surchargé et noirci ; seuls « pass… » et « …ibus » se lisent. Le dernier mot de la phrase, à la fin de la ligne, n'est pas lu (« sirma » ? « firma » ?).}

Correlatarum igitur simpliciorum æquationum tres sunt differentiæ.

prima est ancipitum, in quarum utraq; potestas negatur de homogeneo adfectionis.

altera est contradicentium, in quarum una potestas adficitur affirmate in altera negate.

tertia est \uncertain{inversarum} in quarum una potestas adficitur mul\add{cta} homogeneæ adfectionis in altera e contra \uncertain{homogeneum} adfectionis multatur potestate.
\note{« inversarum » : lecture des lettres (« Jnuersarum »), là où l'on attendrait le nom de la troisième espèce de corrélation ; le même mot revient à la première ligne du feuillet « 21 » (vue 47). « mul… » est au bout de la ligne, au pli, et le reste ne se voit pas. Le bas du feuillet est blanc sous ces lignes.}

\page{47}
\note{Recto. En haut à droite, à l'encre, « 21 ». Le feuillet est plus haut que les deux précédents : ses premières lignes dépassaient déjà au-dessus d'eux aux vues 43 et 45. Sous son bord inférieur dépasse le bas d'un feuillet plus long, avec un signe « ǂ » à gauche et « = A grad… » au milieu ; la même bande reparaît sous le feuillet « 22 » à la vue 49 : c'est le bas du feuillet « 23 » (vue 51), où il est transcrit. À gauche, au pli, des fins de lignes du verso précédent (vue 46). La page continue sans rupture le chapitre XVI « De Syncrisi » ouvert au verso précédent.}

Sive autem ancipitum, sive contradicentium, sive \uncertain{inversarum}, æquationum constitutio, Syncrisi cognoscetur probe, \uncertain{Iucunda} siquidem syncriseos hæc est ratio. quoniam enim quæ uni æquantur inter se æqualia sunt, proponuntur autem potestates duæ adficientes adfectione, eadem datæ omnino magnitudini homogeneæ comparari, ipsum autem adfectum adficiendisve homogeneum utrique potestati iunctum, fieri sub parodicis iisdem gradibus, eademq; parabola. Ergo potestas radicis una cum homogeneo sub suo gradu parodico æquabitur potestati alterius, una cum homogeneo sub suo quoq; gradu parodico, et translatis ad unam ita constitutæ quoq; æquationis partem potestatibus, ad alteram vero iis quæ sub gradibus fiunt adfectionis homogeneis erit rursum æqualitas unde cum differentia aut aggregatum potestatum applicabitur ad differentiam vel aggregatum laterum (ut re vero demum opus indicabit) facit oriri magnitudinem parabolæ æqualem, cuius ideo eadem erit constitutio.
\note{« inversarum » : même mot, et même lecture des lettres, qu'à la fin de la vue 46. « Iucunda » : le mot commence par un grand jambage (I ou T) suivi de « u », puis d'une lettre qui peut être un c ou un e ; lecture douteuse.}

Quod si postea utriusq; positarum æquationum Syncrisi expositarum non iam sub ipsomet parabola exhibeatur sed sub æquo parabolæ (prout eius constitutio exquisita fuit) valore: quod ordinabitur erit consequenter, dato comparationis homogeneo æquale, cuius ideo constitutio non poterit ignorari quoniam oritur illud ex facti ab aggregato vel differentia radicum aut graduum in planum sub radicibus applicatione ad eundem cuius antecedens applicatio facta est terminum, sive differentiam sive aggregatum.

\marginal{\uncertain{Secundum}}
\uncertain{Sit inde} quæ proferuntur parabolæ adfectionum in ancipitibus ex applicatione differentiæ potestatum ad differentiam eorum quos metitur parabola graduum.
\note{Les deux premiers mots de l'alinéa sont soulignés, et la marge de gauche porte en regard, soulignés aussi et suivis d'une croix, un mot lu « Secundum » : sans doute une correction proposée pour les mots soulignés, qui ne sont pas biffés. Les deux leçons sont données, l'une et l'autre douteuses. Devant « ex applicatione », en tête de ligne, un petit trait.}

homogenea vero comparationum ex applicatione differentiæ factorum reciprocè a potestate cuiusq; radicis in gradum alterius ad prædictam cuius altera applicatio facta est graduum differentiam.

Id autem autem obiter est demonstrandum.
\note{« autem » est écrit deux fois, à la suite.}

\subsection*{Problema 1.}

Duarum ancipitum æquationum constitutionem ex Syncrisi dignoscere.

\[ \text{Proponatur } \left.\begin{array}{l} B \text{ parabola in } A \text{ gradum} \\ - A \text{ potestate} \end{array}\right\} \text{æquari } Z \text{ homogeneæ.} \]
\note{Fin de la page ; le bas du feuillet est blanc.}

\page{48}
\note{Verso du feuillet « 21 » ; aucun chiffre en tête. Ce feuillet, le plus haut de la série, ne laisse rien dépasser au-dessus de lui ; à droite de la vue, au pli, les bords d'autres feuillets, avec deux « B » isolés. Dans la marge de gauche transparaît, à l'envers, l'écriture du recto. La page continue sans rupture le problème I laissé au bas du recto (vue 47).}

\[ \text{Et rursus eadem } \left.\begin{array}{l} B \text{ parabolam in } E \text{ gradum} \\ - E \text{ potestate} \end{array}\right\} \text{æquari } Z \text{ homogeneæ eidem} \]

unde par \ill{} gradus et par potestas. oportet ex Syncrisi earum æquationum constitutionem agnoscere.
\note{Le mot après « par » se lit lettre à lettre « de-s-t-i-t » (ou « destit », avec un s long) ; il n'est pas résolu.}

quoniam igitur quæ uni æquantur æqualia sunt inter se manifestum est
\[ \left.\begin{array}{l} B \text{ parabolam in } A \text{ gradum} \\ - A \text{ potestate} \end{array}\right\} \text{æquari} \left\{\begin{array}{l} B \text{ parabolæ in } E \text{ gradum} \\ - E \text{ potestate.} \end{array}\right. \]

Et per antithesin statuendo si placet $A$ maiorem quam $E$ quod quidem hic liberum est, propter præsuppositam radicum ἀμφιβολίαν
\[ \left.\begin{array}{l} A \text{ potestas} \\ - E \text{ potestate} \end{array}\right\} \text{æquabitur} \left\{\begin{array}{l} B \text{ parabolæ in } A \text{ gradum} \\ - B \text{ parabolæ in } E \text{ gradum.} \end{array}\right. \]
\note{« ἀμφιβολίαν » est écrit en lettres grecques, avec son accent.}

Et omnibus per $A$ gradum $-$ $E$ gradu divisis
\[ \left.\begin{array}{c} \begin{array}{l} A \text{ potestas} \\ - E \text{ potestate} \end{array} \\ \hline \begin{array}{l} A \text{ gradum} \\ - E \text{ gradu} \end{array} \end{array}\right\} \text{æquabit. } B \text{ parabolæ.} \]

oritur ergo $B$ parabola ex applicatione differentiæ potestatum ad differentiam graduum ut est constitutum.

porro cum $B$ parabola in $A$ gradum minus $A$ potestate æquetur $Z$ homogeneæ, in locum $B$ parabolæ, subrogetur iam agnitus eius valor, et ea subrogatione æquatio reformetur ergo.
\[ \left.\begin{array}{c} \begin{array}{l} A \text{ potestas in } E \text{ gradum} \\ - E \text{ potestate in } A \text{ gradum} \end{array} \\ \hline \begin{array}{l} A \text{ gradu} \\ - E \text{ gradu} \end{array} \end{array}\right\} \text{æquabitur } Z \text{ homogeneæ.} \]

oritur ergo $Z$ homogenea ex applicatione differentiæ factorum reciprocè a potestate cuiusq; radicis in gradum alterius, ad differentiam graduum, ut est quoq; constitutum.

\subsection*{In Specie.}

\[ \text{Proponatur } \left.\begin{array}{l} B \text{ in } A \\ - A \text{ quadrato} \end{array}\right\} \text{æquari } Z \text{ plano} \]

\[ \text{et rursus } \left.\begin{array}{l} B \text{ in } E \\ - E \text{ quadrato} \end{array}\right\} \text{æquari } Z \text{ plano.} \]

oportet ex Syncrisi earum æquationum constitutionem dignoscere

quoniam igitur quæ uni æquantur æqualia sunt inter se. manifestum est.
\note{Fin de la page ; le bas du feuillet est blanc sous ces lignes. La phrase continue sur le recto suivant (vue 49).}

\page{49}
\note{Recto. En haut à droite, à l'encre, « 22 ». Sous le bord inférieur du feuillet dépasse le bas d'un feuillet plus long, avec, comme à la vue 47, un signe « ǂ » à gauche et « = A grad… » au milieu : c'est le bas du feuillet « 23 » (vue 51), où il est transcrit. À droite de la vue, au pli, des fins de lignes d'autres feuillets. La page continue sans rupture la phrase laissée au bas du verso précédent (vue 48, « manifestum est »).}

\[ \left.\begin{array}{l} B \text{ in } A \\ - A \text{ quadrato} \end{array}\right\} \text{æquari} \left\{\begin{array}{l} B \text{ in } E \\ - E \text{ quadrato} \end{array}\right. \]

Et per antithesin statuendo si placet $A$ maiorem quam $E$ quod quidem hic liberum est propter suppositam radicum ἀμφιβολίαν
\[ \left.\begin{array}{l} A \text{ quadrat.} \\ - E \text{ quadrato} \end{array}\right\} \text{æquabitur} \left\{\begin{array}{l} B \text{ in } A \\ - B \text{ in } E \end{array}\right. \]

Et omnibus per $A - E$ divisis
\[ \left.\begin{array}{c} \begin{array}{l} A \text{ quadratum} \\ - E \text{ quadrato} \end{array} \\ \hline A - E \end{array}\right\} \text{id est } A + E \text{ æquabitur } B. \]

\marginal{quæritur}
Est igitur $B$ summa duorum de quibus oritur laterum \uncertain{oriunda} ex applicatione differentiæ quadratorum ad differentiam laterum, ut est generaliter constitutum.
\note{« oritur » est souligné, et la marge de gauche porte en regard, souligné aussi, « quæritur » : une correction, sans rature du texte ; les deux leçons sont données. « oriunda » : la fin du mot est noyée dans une tache.}

porro cum $\left.\begin{array}{l} B \text{ in } A \\ - A \text{ quadrato} \end{array}\right\}$ æquetur $Z$ plano

si in locum $B$ subrogetur iam agnitus ipsius valor nempe
\[ \begin{array}{c} \begin{array}{l} A \text{ quadratum} \\ - E \text{ quadrato} \end{array} \\ \hline A - E \end{array} \text{ seu } A + E \qquad \begin{array}{c} \begin{array}{l} A \text{ quadrat. in } E \\ - E \text{ quadrato in } A \end{array} \\ \hline A - E \end{array} \text{ id est } E \text{ in } A \]
æquabitur $Z$ plano. est igitur $Z$ planum id quod fit sub duobus de quibus quæritur lateribus, ortum ex differentiæ factorum reciprocè a quadrato unius in radicem alterius applicatione ad differentiam ipsorum laterum ut est quoq; generaliter constitutum.
\note{Le E de « A quadrat. in E » est écrit sur un A.}

\subsection*{Aliud.}

\[ \text{Proponatur } \left.\begin{array}{l} B \text{ plano in } A \\ - A \text{ cubo} \end{array}\right\} \text{æquari } Z \text{ solido} \]

\[ \text{et rursus } \left.\begin{array}{l} B \text{ plano in } E \\ - E \text{ cubo} \end{array}\right\} \text{æquari } Z \text{ solido} \]

oportet ex Syncrisi earum æquationum constitutionem dignoscere quoniam igitur quæ uni æquantur æqualia sunt inter se manifestum est
\[ \left.\begin{array}{l} B \text{ plano in } A \\ - A \text{ cubo} \end{array}\right\} \text{æquari} \left\{\begin{array}{l} B \text{ plano in } E \\ - E \text{ cubo} \end{array}\right. \]

et per antithesin statuendo si placet $A$ maiorem quam $E$ quod quidem hic liberum est propter suppositam radicum ἀμφιβολίαν
\[ \left.\begin{array}{l} A \text{ cubus} \\ - E \text{ cubo} \end{array}\right\} \text{æquabitur} \left\{\begin{array}{l} D \text{ plano in } A \\ - D \text{ plano in } E. \end{array}\right. \]
\note{« D plano », deux fois, là où les lignes précédentes écrivent « B plano » ; une petite croix est tracée au-dessus du premier D et une autre sous le second. Dans la marge intérieure, à gauche, trois « B » isolés de la même encre, l'un en regard de la ligne « manifestum est B plano in A… », les deux autres en regard de ces deux lignes : sans doute des corrections signalées en marge, comme « S/ » à la vue 41 ; le texte n'est pas corrigé ici. Le bas du feuillet est blanc sous ces lignes.}

\page{50}
\note{Verso du feuillet « 22 » ; aucun chiffre en tête. La page continue sans rupture l'exemple « Aliud » laissé au bas du recto (vue 49). Comme au recto, la marge intérieure (ici à gauche) porte des « B » isolés, en regard des lignes où l'écrivain a écrit « D plano », et de petites croix sont tracées au-dessus des D ; voir la note de la vue 49. Dans la marge de gauche transparaît aussi, à l'envers, l'écriture du recto.}

et omnibus per $A$ divisis
\[ \left.\begin{array}{c} \begin{array}{l} A \text{ cubus} \\ - E \text{ cubo} \end{array} \\ \hline A - E \end{array}\right. \text{ id est } \left.\begin{array}{l} A \text{ quadrat.} \\ + E \text{ quadrat.} \\ + A \text{ in } E \end{array}\right\} \text{æquabitur } D \text{ plano.} \]
\note{« per A divisis » : ainsi écrit, là où le diviseur écrit dessous est « A $-$ E ». Le « $+$ » devant « A in E » est tracé sous la ligne et paraît ajouté ; une petite croix au-dessus du D. En marge, en regard, « B ».}

Est igitur $B$ planum aggregatum quadratorum a duobus de quibus quæritur lateribus adiuncto ei quod sub iis fit plano, et oritur ex applicatione differentiæ cuborum ad differentiam laterum ut est generaliter constitutum.

porro cum $\left.\begin{array}{l} D \text{ plano in } A \\ - A \text{ cubo} \end{array}\right\}$ æquetur $Z$ solido.

si in locum $D$ plani subrogetur eius valor nempe
\[ \begin{array}{c} \begin{array}{l} A \text{ cubus} \\ - E \text{ cubo} \end{array} \\ \hline A - E \end{array} \text{ seu } \left|\begin{array}{l} A \text{ quadratum} \\ E \text{ quadratum} \\ A \text{ in } E \end{array}\right. \qquad \begin{array}{c} \begin{array}{l} A \text{ cubus in } E \\ - E \text{ cubo in } A \end{array} \\ \hline A - E \end{array} \text{ id est } \left\{\begin{array}{l} A \text{ quadrat. in } E \\ + E \text{ quadrato in } A \end{array}\right. \]
æquabit $Z$ solido
\note{Dans « seu A quadratum / E quadratum / A in E », les trois termes sont écrits sans signe, joints à gauche par un trait vertical ; dans la marge, en regard des deux derniers, deux croix « $+$ », qui paraissent suppléer les signes. Les deux D de « D plano in A » et de « D plani » portent une petite croix au-dessus, et la marge un « B » en regard de chaque ligne.}

Est igitur $Z$ solidum quod fit ab aggregato laterum in planum sub lateribus. et oritur ex differentiæ ipsorum factorum reciprocè a cubo unius lateris in latus alterius applicatione ad differentiam ipsorum laterum. ut est quoq; constitutum.

Constitutio igitur æquationum propositarum tandem ex Syncrisi agnita est ut faciendum erat.

Placeat exhibere formulam æquationis ancipitis quadraticæ de radice $F$ et $G$. hac minore illa maiore explicabilis.

\[ \text{dixero } \left.\begin{array}{l} F + G \text{ in } A \\ - A \text{ quadrato} \end{array}\right\} \text{æquari } F \text{ in } G \text{ et fieri } A, F \text{ vel } G. \]

Sit $F$ 10 $G$ 2. $A$ 1N. formula æquationis erit 12N $-$ 1Q æquabit 20 explicabilis de 10 vel 2.
\marginal{i}
\note{En marge de gauche, en regard, une petite lettre isolée « i ».}

Placeat exhibere formulam æquationis cubi negati de homogeneo adfectionis sub latere ut sit radix de $F$ vel $G$ explicabilis

\[ \text{dixero } \left.\begin{array}{l} \left.\begin{array}{l} F \text{ quadratum} \\ + G \text{ quadratum} \\ + F \text{ in } G \end{array}\right\} \text{in } A \\ - A \text{ cubo} \end{array}\right\} \text{æquari} \left\{\begin{array}{l} F \text{ in } G \text{ quadratum} \\ + G \text{ in } F \text{ quadratum} \end{array}\right. \]

et fieri $A$ $F$ vel $G$. Sit $F$ 10. $G$ 2. $A$ 1N. formula æquationis erit 124N. $-$ 1C. æquabitur 240. explicabilis de 10 vel 2.
\note{Le dernier « 2 » est repassé. $100 + 4 + 20 = 124$ et $40 + 200 = 240$.}

In contradicentibus coefficientes subgraduales proferuntur ex applicatione differentiæ potestatum ad aggregatum eorum quos sustinent coefficientes graduum.
\note{Fin du texte ; le bas du feuillet est blanc.}

\page{51}
\note{Recto. En haut à droite du feuillet, à l'encre, « 23 ». Au-dessus de son bord supérieur dépasse le haut d'un feuillet plus grand, chiffré « 25 », avec trois lignes (« Ceterum dum aggregatum potestatum applicatur ad differentiam radicum nulla inde orietur magnitudo… ») et, dans sa marge, « a qua » : c'est la tête de la page de la vue 55, où elles sont transcrites. Le feuillet « 23 » est lui-même plus long que les feuillets « 21 » et « 22 » : c'est son bas de page (le signe « ǂ » en marge et les « = A gradu ») qui dépasse sous eux aux vues 47 et 49. La page continue sans rupture le verso précédent (vue 50), qui finissait sur les coefficients « in contradicentibus ».}

homogenea vero datæ nascuntur ex applicatione aggregati factorum reciprocè a potestate radicis unius in gradum alterius ad prædictum cuius altera applicatio facta est graduum aggregatum.
\note{Un trait long traverse « graduum » : c'est la barre du t de « est » prolongée, non une rature.}

Id autem quoq; obiter est demonstrandum.

\subsection*{Problema II}

Duarum contradicentium æquationum constitutionem ex Syncrisi agnoscere

\[ \text{Proponatur } \left.\begin{array}{l} A \text{ potestas} \\ + B \text{ coefficiente in } A \text{ gradum} \end{array}\right\} \text{æquari } Z \text{ homogeneo} \]

\[ \text{et rursus } \left.\begin{array}{l} E \text{ potestas} \\ - \text{eadem } B \text{ coefficiente in } E \text{ gradum} \end{array}\right\} \text{æquari eidem } Z \text{ homogeneo} \]
\note{« eadem » est surchargé d'une tache, comme « earum » à la ligne suivante.}

unde par \ill{} gradus et par potestas. oportet ex Syncrisi earum æquationum constitutionem agnoscere

quoniam igitur quæ uni æquantur æqualia sunt inter se manifestum est
\note{Le mot après « par » est le même qu'à la vue 48, et se lit de même « de-s-t-i-t » ; il reste non résolu.}

\[ \left.\begin{array}{l} A \text{ potestatem} \\ + B \text{ coefficiente in } A \text{ gradum} \end{array}\right\} \text{æquari} \left\{\begin{array}{l} E \text{ potestati} \\ - B \text{ coefficiente in } E \text{ gradu} \end{array}\right. \]

et per antithesin
\[ \left.\begin{array}{l} E \text{ potestatem} \\ - A \text{ potestate} \end{array}\right\} \text{æquari} \left\{\begin{array}{l} B \text{ coefficienti in } E \text{ gradum} \\ + B \text{ coefficiente in } A \text{ gradum.} \end{array}\right. \]

itaq; omnibus per $E$ gradum, plus $A$ gradu divisis
\[ \left.\begin{array}{c} \begin{array}{l} E \text{ potestas} \\ - A \text{ potestate} \end{array} \\ \hline \begin{array}{l} E \text{ gradu} \\ + A \text{ gradu} \end{array} \end{array}\right\} \text{æquabit } B \text{ coefficienti} \]
\note{Devant le « $+$ » du diviseur, un mot ou un signe biffé d'une tache.}

oritur ergo $B$ coefficiens subgradualis ex applicatione differentiæ potestatum ad aggregatum eorum quos sustinet coefficiens graduum ut est constitutum.

\[ \text{porro cum } \left.\begin{array}{l} A \text{ potestas} \\ + B \text{ coefficiente in } A \text{ gradum} \end{array}\right\} \text{æquet } Z \text{ homogeneo.} \]

in locum $B$ coefficientis subrogetur iam agnitus eius valor videlicet
\[ \begin{array}{c} \begin{array}{l} E \text{ potestas} \\ - A \text{ potestate} \end{array} \\ \hline \begin{array}{l} E \text{ gradu} \\ = A \text{ gradu} \end{array} \end{array} \quad \text{et ea subrogatione æquatio reformetur} \]

\[ \text{ergo } A \text{ potestas plus } \begin{array}{c} \begin{array}{l} E \text{ potestas} \\ - A \text{ potestate} \end{array} \\ \hline \begin{array}{l} E \text{ gradu} \\ = A \text{ gradu} \end{array} \end{array} \text{in } A \text{ gradum} \;\Big\}\; \text{æquabitur } Z \text{ homogeneo.} \]
\note{Dans les deux diviseurs du bas de la page, le signe devant « A gradu » est un double trait « = », là où la division précédente écrit « $+$ » ; il est gardé tel qu'il est écrit. En regard, dans la marge de gauche, un « $+$ » puis un signe « ǂ » (deux traits horizontaux barrés d'un trait vertical) : la correction du signe paraît signalée en marge. « homog » est biffé à la fin de la ligne, et « homogeneo » récrit dessous. Le bas du feuillet est blanc sous ces lignes.}

\page{52}
\note{Verso du feuillet « 23 » ; aucun chiffre en tête. Au-dessus du bord supérieur dépasse le haut du verso de la vue 50 (« Et omnibus per A divisis », « A cubus $-$ E cubo id est A quadrat… æquabitur D plano »), transcrit à cette vue ; le B isolé de sa marge se voit aussi. À droite, au pli, des débuts de lignes d'une autre page. La page continue sans rupture le problème II laissé au bas du recto (vue 51).}

hoc est omnibus rite ordinatis
\[ \left.\begin{array}{c} \begin{array}{l} A \text{ potestas in } E \text{ gradum} \\ + E \text{ potestate in } A \text{ gradum} \end{array} \\ \hline \begin{array}{l} E \text{ gradu} \\ + A \text{ gradu} \end{array} \end{array}\right\} \text{æquabitur } Z \text{ homogeneo.} \]

oritur ergo $Z$ homogenea ex applicatione aggregati factorum reciprocè a potestate radicis unius in gradum alterius ad prædictum cuius altera applicatio facta est graduum aggregatum ut est secundo loco enunciatum.

In Inversis plane negatis coefficientes subgraduales proferuntur ex applicatione aggregati potestatum ad aggregatum eorum in quos coefficiens est graduum.
\note{« plane negatis » : lecture des lettres, le premier mot douteux (« plane » ou « plene »).}

homogenea vero datæ nascuntur ex applicatione differentiæ factorum reciprocè a potestate radicis unius in gradum alterius ad prædictum cuius altera applicatio facta est graduum in quos coefficiens subgradualis est aggregatum.

Id autem quoq; obiter est demonstrandum.

\subsection*{Problema III}

Duarum Inversarum æquationum constitutionem ex Syncrisi agnoscere

\[ \text{Proponatur } \left.\begin{array}{l} A \text{ potestas} \\ = B \text{ coefficiente in } A \text{ gradum} \end{array}\right\} \text{æquari } Z \text{ homogeneo} \]
\note{Devant « B coefficiente », un double trait « = » suivi d'un « $+$ » serré contre le B ; en marge de gauche, en regard, un trait « — ». Plus bas, dans la même égalité récrite (« manifestum fit… »), le signe est de nouveau un double trait « = » suivi d'un trait long, et la marge porte en regard un signe fait de plusieurs traits croisés. Le calcul qui suit (« per antithesin ») suppose la soustraction. Les signes sont gardés tels qu'ils se lisent ; le transcripteur ne décide pas.}

\[ \text{et rursus } \left.\begin{array}{l} \text{eadem } B \text{ coefficiens in } E \text{ gradum} \\ - E \text{ potestate} \end{array}\right\} \text{æquari eidem } Z \text{ homogeneo} \]
\note{« rursus » est surchargé. Après « E potestate », « in A gradum » est biffé d'un trait.}

unde intelliguntur $A$ et $E$ pares gradus atq; adeo pares potestates oporteat autem æqualitatum harum constitutionem agnoscere

quoniam igitur quæ uni æquantur æqualia sunt inter se manifestum fit ex iis quæ proponuntur
\[ \left.\begin{array}{l} A \text{ potestatem} \\ = B \text{ coefficiente in } A \text{ gradum} \end{array}\right\} \text{æquari} \left\{\begin{array}{l} B \text{ coefficienti in } E \text{ gradum} \\ - E \text{ potestate} \end{array}\right. \]

Et per antithesin
\[ \left.\begin{array}{l} A \text{ potestatem} \\ + E \text{ potestate} \end{array}\right\} \text{æquari} \left\{\begin{array}{l} B \text{ coefficienti in } A \text{ gradum} \\ + B \text{ coefficiente in } E \text{ gradum} \end{array}\right. \]

Itaq; omnibus per $E$ gradum $+$ $A$ gradu divisis
\[ \left.\begin{array}{c} \begin{array}{l} A \text{ potestas} \\ + E \text{ potestate} \end{array} \\ \hline \begin{array}{l} A \text{ gradu} \\ + E \text{ gradu} \end{array} \end{array}\right\} \text{æquabitur } B \text{ coefficienti.} \]
\note{« A gradu » : le A est écrit sur une autre lettre. Le bas du feuillet est blanc sous ces lignes.}

\page{53}
\note{Recto. En haut à droite du feuillet, à l'encre, « 24 ». Au-dessus de son bord supérieur dépasse de nouveau la tête du feuillet « 25 » (« Ceterum dum aggregatum potestatum… », avec « a qua » en marge), transcrite à la vue 55. À gauche et à droite, au pli, des fins et des débuts de lignes d'autres pages. La page continue sans rupture le problème III laissé au bas du verso précédent (vue 52).}

Ipsum autem est quod primo loco enunciatum.

\[ \text{porro cum } \left.\begin{array}{l} A \text{ potestas} \\ = B \text{ coefficiente subgraduali in } A \text{ gradum} \end{array}\right\} \text{æquetur } Z \text{ homogeneo} \]
\note{Le signe devant « B coefficiente » est de nouveau un double trait « = », tiré long ; en marge de gauche, en regard, le signe « ǂ » à traits multiples de la vue 52. À droite, contre le bord, un mot coupé par le pli (« par… »), qui n'est sans doute pas de cette page.}

In locum $B$ coefficientis subrogatur agnitus valor videlicet
\[ \left.\begin{array}{c} \begin{array}{l} A \text{ potestas} \\ + E \text{ potestate} \end{array} \\ \hline \begin{array}{l} A \text{ gradu} \\ + B \text{ gradu} \end{array} \end{array}\right\} \text{igitur } A \text{ potestas} - \left\{\begin{array}{c} \begin{array}{l} A \text{ potestate in } A \text{ gradum} \\ E \text{ potestate in } A \text{ gradum} \end{array} \\ \hline \begin{array}{l} A \text{ gradu} \\ + E \text{ gradu} \end{array} \end{array}\right. \]
\note{« $+$ B gradu » : le B est souligné, et la marge de gauche porte en regard « $-$ E. » : la correction de la lettre, et peut-être du signe, signalée en marge, sans rature dans le texte.}

æquabitur $Z$ homogeneo hoc est omnibus rite ordinatis
\[ \left.\begin{array}{c} \begin{array}{l} A \text{ potestas in } E \text{ gradum} \\ - E \text{ potestate in } A \text{ gradum} \end{array} \\ \hline \begin{array}{l} A \text{ gradu} \\ + E \text{ gradu} \end{array} \end{array}\right\} \text{æquabitur } Z \text{ homogeneo.} \]

Quod ipsum est quod \uncertain{secundo} loco enunciatum.
\note{« secundo » : les lettres se lisent « sit nudo », en deux groupes ; la même formule, à la vue 52, porte « ut est secundo loco enunciatum ».}

In Inversis autem quarum una per affirmationem adficitur, altera per negationem, coefficientes subgraduales proferuntur ex applicatione aggregati potestatum ad differentiam eorum in quos coefficiens est graduum

homogenea vero datæ nascuntur ex applicatione aggregati factorum reciprocè a potestate radicis unius in gradum alterius ad prædictam cuius altera applicatio facta est graduum in quos coefficiens est differentiam.
\note{« aggregati » est traversé d'un trait, qui peut être une rature ; aucun mot n'est écrit à sa place.}

Id autem quoq; obiter est demonstrandum.

\[ \text{Proponatur siquidem } \left.\begin{array}{l} A \text{ potestas} \\ + B \text{ coefficiente in } A \text{ gradum} \end{array}\right\} \text{æquari } Z \text{ homogeneo.} \]

\[ \text{et rursus } \left.\begin{array}{l} \text{eadem } B \text{ coefficiens in } E \text{ gradum} \\ - E \text{ potestate} \end{array}\right\} \text{æquari eidem } Z \text{ homogeneo} \]

unde intelliguntur $A$ et $E$ pares gradus atq; adeo pares potestates oporteat autem æqualitatum harum constitutionem agnoscere.

quoniam igitur quæ uni æquantur æqualia sunt inter se manifestum fit ex iis quæ proponuntur
\[ \left.\begin{array}{l} A \text{ potestatem} \\ - E \text{ potestate} \end{array}\right\} \text{æquari} \left\{\begin{array}{l} B \text{ coefficienti in } E \text{ gradum} \\ - B \text{ coefficienti in } A \text{ gradum.} \end{array}\right. \]
\note{« $-$ E potestate » : un « $+$ » est écrit en marge de gauche, en regard ; la division qui suit porte « $+$ A potestate ».}

Itaq; omnibus per $E$ gradum minus $A$ gradu divisis
\[ \left.\begin{array}{c} \begin{array}{l} E \text{ potestas} \\ + A \text{ potestate} \end{array} \\ \hline \begin{array}{l} E \text{ gradu} \\ - A \text{ gradu} \end{array} \end{array}\right\} \text{æquabitur } B \text{ coefficienti.} \]
\note{Le bas du feuillet est blanc sous ces lignes.}

\page{54}
\note{Verso du feuillet « 24 » ; aucun chiffre en tête. Au-dessus du bord supérieur dépasse de nouveau le haut du verso de la vue 50 (« Et omnibus per A divisis… æquabitur D plano »), transcrit à cette vue. La page continue sans rupture la démonstration laissée au bas du recto (vue 53), puis ferme la matière des corrélations par une remarque générale.}

Quod ipsum est quod primo loco enunciatum.

\[ \text{porro cum } \left.\begin{array}{l} A \text{ potestas} \\ + B \text{ coefficiente in } A \text{ gradum} \end{array}\right\} \text{æquetur } Z \text{ homogeneo} \]

In locum $B$ coefficientis subgradualis subrogetur agnitus valor videlicet
\[ \begin{array}{c} \begin{array}{l} E \text{ potestas} \\ + A \text{ potestate} \end{array} \\ \hline \begin{array}{l} E \text{ gradu} \\ - A \text{ gradu} \end{array} \end{array} \]
\[ \text{igitur } A \text{ potestas plus } \left.\begin{array}{c} \begin{array}{l} E \text{ potestate} \\ + A \text{ potestate} \end{array} \\ \hline \begin{array}{l} E \text{ gradu} \\ - A \text{ gradu} \end{array} \end{array}\right\} \text{in } A \text{ gradum} \;\Big\}\; \text{æquabitur } Z \text{ homogeneo.} \]

hoc est omnibus rite ordinatis
\[ \left.\begin{array}{c} \begin{array}{l} A \text{ potestas in } E \text{ gradum} \\ \cdot\ E \text{ potestate in } A \text{ gradum} \end{array} \\ \hline \begin{array}{l} \cdot\ E \text{ gradu} \\ - A \text{ gradu} \end{array} \end{array}\right\} \text{æquabit } Z \text{ homogeneo.} \]
\note{Devant « E potestate in A gradum » et devant « E gradu », au lieu d'un signe, un simple point ; en marge de gauche, en regard, une croix « $+$ ». Le signe n'est pas suppléé.}

\marginal{Quoniam in ancipitibus differentia potestatum applicata ad differentiam graduum æquatur $B$ parabolæ in quadraticis igitur $\begin{array}{c} \begin{array}{l} A q. \\ - E q. \end{array} \\ \hline A - E \end{array}$ erit $A + E$. in cubis $\begin{array}{c} \begin{array}{l} A \text{ cub.} \\ - E \text{ cub.} \end{array} \\ \hline A - E \end{array}$ erit $\begin{array}{l} A q. \\ + E \text{ in } A \\ + E q. \end{array}$ in quadrato quadratis $\begin{array}{c} \begin{array}{l} A qq. \\ - E qq. \end{array} \\ \hline A - E \end{array}$ \(\}\) $\begin{array}{l} A \text{ cub.} \\ + E \text{ in } A q. \\ + E q. \text{ in } A \\ + E \text{ cub.} \end{array}$ ex capite de genesi potestatum a binomia radice. Et pro cubis vide Cap. \ill{} X. lib. 2.}
\note{Longue note de la marge de gauche, de la même main, rattachée par un renvoi (un signe fait d'un « E » et d'un petit « A », répété devant la ligne « Quod ipsum est ») à la ligne qui suit. Dans le terme des carrés-carrés, un premier essai est biffé (« B … A », « $-$ A qq. ») et « A cub. » est écrit sur « A qq. ». Les abréviations sont gardées : « A q. » (A quadratum), « A qq. » (quadrato-quadratum). La dernière ligne, en partie sous une tache, se lit « Et pro cubis vide … X. lib. 2. » ; le mot qui suit « vide » n'est pas lu, et le « X » est barré d'un trait qui peut être une rature.}

Quod ipsum est quod \uncertain{secundo} loco enunciatum.
\note{Même graphie qu'à la vue 53 (« sit nudo »).}

Et ancipites quidem \uncertain{Divinus} omnino et in omni potestatum ordine efficaces quoniam agnita semel magistra, liberior est ad commentum \uncertain{transitus} ut ex ex \uncertain{retentis} clarum est.
\note{« Divinus » est écrit avec une majuscule, lecture des lettres. La fin de la phrase est douteuse : « transitus » peut se lire « hauſitus » ; « ex ex » est écrit deux fois ; « retentis » (ou « relictis ») est incertain. Un signe en forme de « 7 » barré suit le point final, sans renvoi repéré.}

At contradicentes et Inversæ in \uncertain{lucidum} efficaces sunt in \uncertain{lucidum} minime \uncertain{idq;} e gradu potestatis definitur et potius contradicentes sub latere tandem sunt efficaces si hæreant in quadrato vel potestatibus deinde alternis videlicet quadrato quadrato, cubocubo et deinceps continuo ordine.
\note{« lucidum » (deux fois) : lecture des lettres, douteuse. « idq; » : de même.}

In base autem et potestatibus deinde alternis videlicet cubo quadratocubo et deinceps continuo ordine consimiliter \uncertain{pugnas} \struck{et \ill{}} Inefficaces, quoniam nihil solatii præsidiive afferunt analystæ quo sibi pateat \ill{} facilior felicior aditus.
\note{Après « pugnas », deux mots biffés d'un trait, dont le premier est « et ». Le mot non lu après « pateat » est écrit en lettres grecques (quelque chose comme « ζητ… »), sans que la lecture aboutisse.}

Ut contra Inversæ sub latere cum utraq; negata est \struck{efficaces} sunt, si hæreant sub cubo, quadratocubo.
\note{« efficaces » est biffé d'un trait, à la fin de la ligne.}

Causa est quod quando differentia potestatum applicatur ad differentiam radicum quæ inde oritur magnitudo æqua fit effectis continue proportionalibus a radicibus iisdem secundum affirmata in quocunq; potestatum ordine, secundum ea quæ exposita sunt in Capite de genesi potestatum a binomia radice

At cum differentia potestatum applicatur ad aggregatum radicum quæ inde oritur magnitudo æqua fit effectis continue proportionalibus prosaphæresi alterna tantummodo, cum potestates sunt cubi quadratocubi et deinceps ordine intermittente alterno.
\note{Une longue parenthèse à la plume, dans la marge de gauche, enferme les deux derniers alinéas. En regard, dans la même marge, une seconde note :}
\marginal{Hæc omnia \uncertain{evidentiam} ac demonstrationem habent ex notis in capite de genesi potestatum a binomia radice.}
\note{« evidentiam » : le mot se termine par un « a » surmonté d'un trait (« evidentiā »), lecture douteuse ; « ex notis » : le groupe « ex » est un signe. Le bas du feuillet est blanc sous ces lignes.}

\page{55}
\note{Recto. En haut à droite, à l'encre, « 25 ». Ce feuillet, plus haut que les précédents, laissait dépasser ses premières lignes au-dessus d'eux aux vues 51 et 53. Sous son bord inférieur dépasse le bas d'un feuillet plus long, avec quatre lignes (« …est composita ex tribus proportionalibus… quod fit ductu alterius extremæ in quadratum compositæ e duabus reliquis. et fit A composita e duabus primis vel e duabus postremis. ») ; c'est le bas du feuillet « 27 », transcrit à la vue 59. La page ne reprend pas le verso précédent (vue 54), qui finissait sur une remarque close : elle commence par une phrase de même matière, puis ouvre un nouveau chapitre.}

\marginal{a qua}
Ceterum dum aggregatum potestatum \uncertain{applicatur} ad differentiam radicum nulla inde orietur magnitudo quæ proportionalibus continue effectis sive per gradus sive numero pares sive impares clymactericæ magnitudines \uncertain{prodeant}.
\note{« applicatur » est souligné, et « a qua » est écrit en regard dans la marge de gauche ; « quæ » est souligné de même : une correction de la construction proposée en marge (« a qua » pour « quæ » ?), sans rature dans le texte. « orietur » est surchargé. « clymactericæ » : ainsi écrit. « prodeant » : le mot est surchargé et noirci.}

\subsection*{Syncriticæ doctrinæ Geometrica phrasis. Cap. XVII.}

Hæc autem syncritica Iudicia exornantur et expoliuntur per aliquot analogias, quibus veritatis \struck{geometrica} \uncertain{nitescunt} et evidentior fit \struck{et} paratus.
\note{« geometrica » est biffé d'un trait, que traverse une tache ; « nitescunt » (ou « mitescunt ») est la lecture des lettres. « et », devant « paratus », est biffé ; « paratus » (ou « paratum ») est souligné.}

\subsection*{Prop. 1}

Si fuerint series trium proportionalium

Est ut prima ad tertiam ita quadratum e prima ad quadratum e secunda.

Et si quatuor. Est ut prima ad quartam ita cubus primæ ad cubum e secunda.

Et si quinq;. Est ut prima ad quintam ita quadratoquadratum e prima ad quadratoquadratum e secunda.

Et si sex. Est ut prima ad sextam ita quadratocubus e prima ad quadratocubum e secunda

Et si septem. Est ut prima ad septimam ita cubocubus e prima ad cubocubum e secunda.

Nam quadratum est potestas rationis duplicatæ, cubus triplicatæ, quadratoquadratum quadruplicatæ et ut alibi annotatum est. Itaq; potestatibus singulis sua adscribuntur proportionalium series secundum earum conditionem

\subsection*{Prop. II}

Si fuerint series trium proportionalium.

Est ut prima ad tertiam ita aggregatum quadratorum a singulis duabus primis ad aggregatum quadratorum a singulis duabus postremis

Et si quatuor.

Est ut prima ad quartam ita aggregatum cuborum a singulis tribus primis ad aggregatum cuborum a singulis tribus postremis

Et si quinq;

Est ut prima ad quintam ita aggregatum quadrato-quadratorum a singulis quatuor primis ad aggregatum quadratoquadratorum a singulis quatuor postremis.
\note{Fin du texte du feuillet ; le bas est blanc.}

\marginal{\uncertain{Circa} \quad $\left.\begin{array}{l} B \text{ pl in } A \\ \ldots \end{array}\right\}$ $Z$ s. \quad Sit $A$ cub. \ill{} \quad et similiter $A$ esse $F$ \quad Itaq; ex præcedenti capite $\left.\begin{array}{l} F q. \\ F \text{ in } E \\ E q. \end{array}\right\}$ $B$ plano \struck{et} erit $E$ radix huius æquationis altera \quad Itaq; \uncertain{instituta} \ill{} \quad $\left.\begin{array}{l} B \text{ plano} \\ - F q. \end{array}\right\}$ $\left\{\begin{array}{l} F \text{ in } E \\ + E q. \end{array}\right.$ sunt autem proportionalia plana $F q.$, $F \text{ in } E$, $E q.$ \quad \ill{} \quad æquetur $A q.$ utique ut $F q.$ ad $A q.$ ita $A q.$ ad $\frac{A qq.}{F q.}$ \quad \struck{\ill{} $F q. + A q. + \frac{A qq.}{F q.}$ \(\}\) $B$ pl.} \quad \struck{$F qq.$ $+$ $F q.$ in $A q.$} \(\}\) $B$ p in $F q.$ \quad et $\left.\begin{array}{l} F q. \text{ in } A q. \\ + A qq. \end{array}\right|$ $\begin{array}{l} B \text{ pl in } F q. \\ + F qq. \end{array}$}
\note{La marge de gauche porte, sur la moitié de sa hauteur et de la même main, un calcul très corrigé : plusieurs lignes sont barrées d'un trait oblique, des lettres sont écrites les unes sur les autres, et la lecture n'en est donnée que pour ce qui se lit. Il semble appliquer aux carrés-carrés la syncrise du chapitre précédent, avec des « plans proportionnels » $F q.$, $F$ in $E$, $E q.$ ; le transcripteur ne le reconstitue pas. Le premier mot, souligné, est lu « Circa » avec doute ; sous « Sit A cub. », un « Z » ; en bas de la marge, un trait isolé en forme de « 1 ».}

\page{56}
\note{Verso du feuillet « 25 » ; aucun chiffre en tête. Sous son bord inférieur dépasse le bas du verso de la vue 54 (« Prosaphæresi alterna tantummodo, cum potestates sunt cubi quadratocubi et deinceps ordine intermittente alterno », avec la parenthèse de la marge), transcrit à cette vue. À droite, au pli, les débuts de lignes d'une note de marge d'un autre feuillet (« Sit di… / ex quad… / primam / in qua… / differen… / a mediis »), qui est de la page de la vue 57. La page continue sans rupture la proposition II laissée au bas du recto (vue 55).}

Et si sex.

Est ut prima ad sextam ita aggregatum quadratocuborum a singulis quinq; primis ad aggregatum quadratocuborum a singulis quinq; postremis

Et si septem.

Est ut prima ad septimam ita aggregatum cubocuborum a singulis sex primis ad aggregatum cubocuborum a singulis sex postremis

quoniam enim ut prima ad tertiam ita est quadratum primæ ad quadratum secundæ, ut vero quadratum primæ ad quadratum secundæ ita quadratum secundæ ad quadratum tertiæ, ergo per synæresim est ut prima ad tertiam ita quadratum primæ plus quadrato secundæ ad quadratum secundæ plus quadrato tertiæ. nec dissimiliter in reliquis seriebus proportionalium et conditionariis secundum rationem extremarum potestatibus licet arguere et ratiocinari.

\subsection*{Prop. III}

Si fuerint series trium proportionalium

Est ut prima ad tertiam ita quadratum compositæ e duabus primis ad quadratum compositæ e duabus postremis.
\note{Après « postremis », un petit signe en forme de « c » barré, et un trait dans la marge de droite, sans renvoi repéré.}

Et si quatuor.

Est ut prima ad quartam ita cubus compositæ e tribus primis ad cubum compositæ ex tribus postremis

Et si quinq;

Est ut prima ad quintam ita quadratoquadratum compositæ e quatuor primis ad quadratoquadratum compositæ e quatuor postremis

Et si sex

Est ut prima ad sextam ita quadratocubus compositæ e quinq; primis ad quadratocubum compositæ e quinq; postremis

Et si septem.

Est ut prima ad septimam ita cubocubus compositæ ex sex primis ad cubocubum compositæ ex sex postremis.

Per synæresim enim est ut prima ad secundam ita composita ex omnibus minus prima ad compositam ex omnibus minus alia extrema quotcunq; proportionalium sit series. In unaquaq; autem serie exponitur sua potestas conditionaria id est tantuplicatæ rationis quam extremarum comparatio exigit.

\subsection*{Prop. IIII}

Si fuerit series trium proportionalium est ut prima ad tertiam ita differentia quadratorum a duabus primis ad differentiam quadratorum a duabus postremis.
\note{Fin du texte du feuillet ; le bas est blanc au-dessus de la bande qui dépasse.}

\page{57}
\note{Recto. En haut à droite, à l'encre, « 26 ». Sous son bord inférieur dépasse le bas d'un feuillet plus long, avec quatre lignes (« est et composita ex tribus proportionalibus et solidum quod fit ductu alterius extremæ in quadratum compositæ e duabus reliquis. et fit A composita e duabus primis vel e duabus postremis. ») ; la même bande dépasse sous le feuillet « 25 » à la vue 55. C'est le bas du feuillet « 27 », transcrit à la vue 59. À droite, au pli, des fragments d'un autre feuillet (« …tiona », « …vel 10 », « 10. »). La page continue sans rupture la proposition IIII laissée au bas du verso précédent (vue 56).}

Et si quatuor est ut prima ad quartam ita differentia cuborum a tribus primis alterne sumptis ad differentiam cuborum a tribus postremis alterne sumptis

Et si quinq; est ut prima ad quintam ita differentia quadratoquadratorum a quatuor primis alterne sumptis ad differentiam quadratoquadratorum a quatuor postremis alterne sumptis.

Et si sex est ut prima ad sextam ita differentia quadratocuborum a quinq; primis alterne sumptis ad differentiam quadratocuborum a quinq; postremis alterne sumptis.

Et si septem.

Est ut prima ad septimam ita differentia cubocuborum a sex primis alterne sumptis ad differentiam cubocuborum a sex postremis alterne sumptis.

\subsection*{Prop. V.}

Si fuerit series trium proportionalium est ut prima ad tertiam ita quadratum differentiæ duarum primarum ad quadratum differentiæ duarum postremarum.

Et si quatuor

Est ut prima ad quartam ita cubus differentiæ trium primarum alterne sumptarum ad cubum differentiæ trium postremarum alterne sumptarum.

Et si quinq;

Est ut prima ad quintam ita quadratoquadratum differentiæ quatuor primarum alterne sumptarum ad quadratoquadratum differentiæ quatuor postremarum alterne sumptarum.

Et si sex

Est ut prima ad sextam ita quadratocubus differentiæ quinq; primarum alterne sumptarum ad quadratocubum differentiæ quinq; postremarum alterne sumptarum.

Et si septem.

Est ut prima ad septimam ita cubocubus differentiæ sex primarum alterne sumptarum ad cubocubum differentiæ sex postremarum alterne sumptarum.

\subsection*{Prop. VI}

\marginal{Sit differentia solidorum ex quadrato quartæ in primam et quadrato primæ in quartam æqualis differentiæ cuborum a mediis}
Si fuerint quatuor continue proportionales solidum compositum ex cubo quartæ et triplo cubo secundæ differt a solido composito ex cubo primæ et triplo cubo tertiæ per cubum differentiæ extremarum
\note{La note de la marge de gauche, de la même main, est écrite en regard de l'énoncé et paraît en proposer une variante ou un complément ; elle n'est rattachée au texte par aucun signe. Le « Si » initial de l'énoncé est caché sous la fin de la note. Fin du texte du feuillet.}

\page{58}
\note{Verso du feuillet « 26 » ; aucun chiffre en tête. Sous son bord inférieur dépasse de nouveau le bas du verso de la vue 54 (« Prosaphæresi alterna tantummodo… ordine intermittente alterno »), transcrit à cette vue. La page donne la démonstration de la proposition VI énoncée au bas du recto (vue 57), puis ouvre le chapitre XVIII.}

Sint quatuor continue proportionales $B$, $D$, $F$, $G$. et sit $G$ maior extrema dico
\[ \left.\begin{array}{l} G \text{ cubum} \\ + D \text{ cubo ter} \\ - B \text{ cubo} \\ - F \text{ cubo ter} \end{array}\right\} \text{æquari } G - B \text{ cubo} \]

Cum vero $G - B$ cubus constat
\[ \left.\begin{array}{l} G \text{ cubo} \\ - B \text{ in } G \text{ quadrato ter} \\ + B \text{ quadrato in } G \text{ ter} \\ - B \text{ cubo} \end{array}\right\} \text{comparatur autem} \quad \begin{array}{l} G \text{ cubo} \\ + D \text{ cubo ter} \\ - B \text{ cubo} \\ - F \text{ cubo ter} \end{array} \]

superest ut
\[ \left.\begin{array}{l} B \text{ quadratum in } G \text{ ter} \\ - B \text{ in } G \text{ quadratum ter} \end{array}\right\} \text{æquetur} \left\{\begin{array}{l} D \text{ cubo ter} \\ - F \text{ cubo ter.} \end{array}\right. \]

Id autem ita se habet nam $B$ quadratum in $G$ est $D$ cubus et $B$ in $G$ quadratum $F$ cubus

\subsection*{Prop. VII}

Si fuerint quatuor continue proportionales. Solidum compositum ex cubo quartæ et triplo cubo secundæ \uncertain{adiunctum} solido composito ex cubo primæ et triplo cubo tertiæ æquatur cubo aggregati extremarum.

eadem enim viget quæ in antecedente theoremate demonstratio.
\note{« adiunctum » : les lettres se lisent « admertum » ; le sens (la somme des deux solides est le cube de la somme des extrêmes) demande « adiunctum ».}

\subsection*{Æquationum ancipitum constitutiva Cap. XVIII}

agnitæ per syncrisin ancipites æquationes, aut plasmatis resolutione, aut denique \uncertain{catholes} vi isto modo sese habent
\note{« catholes » : lecture des lettres, douteuse (« catholicæ » ?).}

\subsection*{De affectis sub depressiore gradu seu latere.}

\subsection*{Prop. I}

\[ \text{Si } \left.\begin{array}{l} B \text{ in } A \\ - A \text{ quadrato} \end{array}\right\} \text{æquetur } Z \text{ plano.} \]

Est $B$ composita e duabus lateribus sub quibus quod fit æquum est $Z$ plano, et fit $A$ latus maius minusve.

Sunto duo latera 2 et 10 dicetur 12N $-$ 1Q æquari 20 et fiet 1N. 2. vel 10.
\note{Le Z de « Z plano », dans la deuxième phrase, est repassé et barré ; la marge de gauche porte en regard un « Z » : la correction est signalée en marge.}

\subsection*{Prop. II}

\[ \text{Si } \left.\begin{array}{l} B \text{ planum in } A \\ - A \text{ cubo} \end{array}\right\} \text{æquetur } Z \text{ solido} \]
\note{Le « A » de « in A » est récrit plus grand, et un trait est tiré devant l'accolade.}

Est $B$ plano compositum ex quadratis trium proportionalium. Et $Z$ solidum quod fit ductu alterius extremæ in aggregatum quadratorum a reliquis.

Et fit $A$ prima vel tertia

Sunto proportionales. 2. $L$ 20. 10. dicetur 124N $-$ 1C æquari 240 et fiet 1N. 2 vel 10.
\note{« $L$ 20 » : le signe en forme de L, déjà rencontré dans la marge de la vue 42, note ici la racine : $2$, $\sqrt{20}$, $10$ sont en proportion continue ; $4 + 20 + 100 = 124$ et $2 \cdot (20 + 100) = 240$. Fin du texte de la page.}

\page{59}
\note{Recto. À droite, à l'encre, « 27 », à la hauteur du titre « Prop. III ». Au-dessus du bord supérieur de ce feuillet dépasse la tête d'un feuillet plus grand, chiffré « 29 » : « Prop. XV. / Si B plano in A cubum [$-$] A quadratocubo \(\}\) æquetur Z planosolido », la seconde ligne coupée par le bord ; ce feuillet n'est pas dans les vues de ce lot (voir l'en-tête), et ces lignes ne sont pas transcrites ici comme texte de la page. La vue est plus large que les autres : à gauche, le bord perforé d'un film. L'encre de cette page est plus pâle que celle des feuillets précédents ; la main est la même. La page continue sans rupture le chapitre XVIII (vue 58, prop. II) ; c'est son bas de page, plus long que les feuillets « 25 » et « 26 », qui dépasse sous eux aux vues 55 et 57.}

\subsection*{Prop. III}

\[ \text{Si } \left.\begin{array}{l} B \text{ solidum in } A \\ - A \text{ quadratoquadrato} \end{array}\right\} \text{æquetur } Z \text{ planoplano} \]

Est $B$ solidum compositum ex cubis quatuor continue proportionalium Et $Z$ planoplanum quod fit ductu alterius extremæ in aggregatum cuborum a tribus reliquis. Et fit $A$ prima vel quarta.

Sunto continue proportionales 2. $L$C 40. $L$C 200. 10 dicetur 1248N $-$ 1QQ æquari 2480 et fiet 1N. 2 vel 10.
\note{« $L$C 40 » : racine cubique de 40. $8 + 40 + 200 + 1000 = 1248$ et $2 \cdot (40 + 200 + 1000) = 2480$.}

\subsection*{Prop. IIII}

\[ \text{Si } \left.\begin{array}{l} B \text{ planoplanum in } A \\ - A \text{ quadratocubo} \end{array}\right\} \text{æquetur } Z \text{ planosolido} \]

Est $B$ planoplanum compositum ex quadratoquadratis quinq; continue proportionalium et $Z$ planosolidum quod fit ductu alterius extremæ in aggregatum quadratoquadratorum a quatuor reliquis, et fit $A$ prima vel quinta

Sunto continue proportionales.

\[ \begin{array}{ccccc} \text{I} & \text{II} & \text{III} & \text{IIII} & \text{V} \\ 2. & L\text{QQ } 80. & L\,20. & L\text{QQ } 2000. & 10. \end{array} \]

dicetur 12496N. $-$ 1QC æquari 24960 et fiet 1N. 2. vel 10.
\note{Les rangs I à V sont écrits au-dessus des termes. « $L$QQ » : racine carrée-carrée. $16 + 80 + 400 + 2000 + 10000 = 12496$ et $2 \cdot (80 + 400 + 2000 + 10000) = 24960$.}

\subsection*{Prop. V.}

\[ \text{Si } \left.\begin{array}{l} B \text{ planosolidum in } A \\ - A \text{ cubocubo} \end{array}\right\} \text{æquetur } Z \text{ solidosolido.} \]

Est $B$ planosolidum compositum ex quadratocubis sex continue proportionalium et $Z$ solidosolidum quod fit ductu alterius extremæ in aggregatum quadratocuborum e quinq; reliquis et fit $A$ prima vel sexta.

Sunto continue proportionales

\[ \begin{array}{cccccc} \text{I} & \text{II} & \text{III} & \text{IIII} & \text{V} & \text{VI} \\ 2. & L\text{QC } 160. & L\text{QC } 800. & L\text{QC } 4000. & L\text{QC } 20{,}000. & 10 \end{array} \]

dicetur 124,992N. $-$ 1CC. æquari 124,960. et fiet 1N. 2. vel 10.

\marginal{$L$QC. 160. \quad 249920}
\note{Le premier chiffre du second terme est surchargé (« 360 » ou « 160 »), et la marge de gauche récrit en regard « $L$QC. 160. » ; le terme est bien $\sqrt[5]{160}$. « 124,960 » est souligné, et la marge porte en regard, sous un trait, « 249920 » : c'est la valeur juste, $2 \cdot (160 + 800 + 4000 + 20000 + 100000) = 249920$ ; le texte, qui n'est pas corrigé, reste tel qu'il est écrit. $32 + 160 + 800 + 4000 + 20000 + 100000 = 124992$.}

\subsection*{De Adfectis sub elatiore gradu seu latere recipro.}
\note{« recipro » : ainsi écrit, pour « reciproco ».}

\subsection*{Prop. VI.}

\[ \text{Si } \left.\begin{array}{l} B \text{ in } A \text{ quadratum} \\ - A \text{ cubo} \end{array}\right\} \text{æquetur } Z \text{ solido} \]

est $B$ composita ex tribus proportionalibus et $Z$ solidum quod fit ductu alterius extremæ in quadratum compositæ e duabus reliquis. et fit $A$ composita e duabus primis vel e duabus postremis.
\note{Fin du texte du feuillet.}

\page{60}
\note{Verso du feuillet « 27 » ; aucun chiffre en tête. Le haut de la vue, au-dessus du bord du feuillet, est vide. À droite, au pli, les débuts de lignes d'une note de marge d'un autre feuillet (« x id est A… / compositæ… / secundæ… / et tertiæ… », et un « 5 » souligné), qui n'est pas de cette page. La page continue sans rupture la proposition VI laissée au bas du recto (vue 59), dont elle donne l'exemple.}

Sint proportionales 1. 2. 4. dicetur 7Q $-$ 1C æquari 36 et fit 1N. 3. vel 6.
\note{Les rangs I, II, III sont écrits au-dessus des trois termes, comme dans les exemples suivants. « 7Q » : le Q a ici la forme d'un 2 (« 72 ») ; avec les racines 3 et 6, $63 - 27 = 36$ et $252 - 216 = 36$.}

\subsection*{Prop. VII}

\[ \text{Si } \left.\begin{array}{l} B \text{ in } A \text{ cubum} \\ - A \text{ quadratoquadrato} \end{array}\right\} \text{æquetur } Z \text{ planoplano.} \]
\note{« B in A cubum » est souligné d'un trait, et un point est posé devant « $-$ A quadratoquadrato ».}

Est $B$ composita e quatuor continue proportionalibus et $Z$ planoplanum quod fit ductu alterius extremæ in cubum compositæ e tribus reliquis et fit $A$ composita e tribus primis vel e tribus postremis.

Sunto continue proportionales 1. 2. 4. 8.

dicetur 15C. $-$ 1QQ. æquari 2744. et fiet 1N. 7 vel 14.

\subsection*{Prop. VIII}

\[ \text{Si } \left.\begin{array}{l} B \text{ in } A \text{ quadratoquadratum.} \\ - A \text{ quadratocubo.} \end{array}\right\} \text{æquetur } Z \text{ planosolido} \]

Est $B$ composita ex quinq; continue proportionalibus et $Z$ planosolidum fit e ductu alterius extremæ in quadratoquadratum compositæ e quatuor reliquis et fit $A$ composita ex quatuor primis vel e quatuor postremis.

Sint continue proportionales 1. 2. 4. 8. 16.

31QQ $-$ 1QC. æquetur 810,000 et fit 1N. 15 vel 30.
\note{« 31QQ » : écrit « 3122 », les deux Q ayant la même forme qu'au premier exemple. $16 \cdot 15^4 = 810000$.}

\subsection*{Prop. IX.}

\[ \text{Si } \left.\begin{array}{l} B \text{ in } A \text{ quadratocubum} \\ - A \text{ cubocubo} \end{array}\right\} \text{æquetur } Z \text{ solidosolido.} \]

Est $B$ composita e sex continue proportionalibus et $Z$ solidosolidum quod fit ductu alterius extremæ in quadratocubum compositæ e reliquis

Et fit $A$ composita e quinq; primis vel e quinq; postremis

Sint continue proportionales. 1. 2. 4. 8. 16. 32.

63QC $-$ 1CC. æquatur 916,132,832.

et fit 1N. 31. vel 62.
\note{$32 \cdot 31^5 = 916\,132\,832$ ; le calcul est juste.}

\subsection*{De affectis sub gradibus intermediis quæ per interpretationem deprimuntur.}

\subsection*{Prop. X.}

\[ \text{Si } \left.\begin{array}{l} B \text{ plano in } A \text{ quadratum.} \\ - A \text{ quadratoquadrato} \end{array}\right\} \text{æquetur } Z \text{ planoplano} \]
\note{Le texte s'arrête ici, au milieu de la page : l'énoncé de la proposition X n'a ni démonstration ni exemple sur ce verso, dont le bas est blanc. La suite, s'il y en a une, est au-delà de ce lot.}

\end{document}
